% =====================================================================
% shared/datos-modelo.tex
% ---------------------------------------------------------------------
% Nomenclatura, ecuaciones y valores de parámetros del modelo de
% anclaje y ajuste del subastador. Fuente única para la memoria
% (docs/MemoriaMHurtado) y la presentación de defensa
% (docs/Presentations/presentation-template), de modo que un valor o una
% ecuación se escriban una sola vez.
%
% Uso:
%   \newcommand{\refcita}[1]{\citeA{#1}}   % memoria (newapa/apacite)
%   % =====================================================================
% shared/datos-modelo.tex
% ---------------------------------------------------------------------
% Nomenclatura, ecuaciones y valores de parámetros del modelo de
% anclaje y ajuste del subastador. Fuente única para la memoria
% (docs/MemoriaMHurtado) y la presentación de defensa
% (docs/Presentations/presentation-template), de modo que un valor o una
% ecuación se escriban una sola vez.
%
% Uso:
%   \newcommand{\refcita}[1]{\citeA{#1}}   % memoria (newapa/apacite)
%   % =====================================================================
% shared/datos-modelo.tex
% ---------------------------------------------------------------------
% Nomenclatura, ecuaciones y valores de parámetros del modelo de
% anclaje y ajuste del subastador. Fuente única para la memoria
% (docs/MemoriaMHurtado) y la presentación de defensa
% (docs/Presentations/presentation-template), de modo que un valor o una
% ecuación se escriban una sola vez.
%
% Uso:
%   \newcommand{\refcita}[1]{\citeA{#1}}   % memoria (newapa/apacite)
%   % =====================================================================
% shared/datos-modelo.tex
% ---------------------------------------------------------------------
% Nomenclatura, ecuaciones y valores de parámetros del modelo de
% anclaje y ajuste del subastador. Fuente única para la memoria
% (docs/MemoriaMHurtado) y la presentación de defensa
% (docs/Presentations/presentation-template), de modo que un valor o una
% ecuación se escriban una sola vez.
%
% Uso:
%   \newcommand{\refcita}[1]{\citeA{#1}}   % memoria (newapa/apacite)
%   \input{../../shared/datos-modelo}
% o
%   \newcommand{\refcita}[1]{\citet{#1}}   % presentación (natbib)
%   \input{../../../shared/datos-modelo}
%
% Los macros de tabla entregan sólo el `tabular`, sin flotante ni
% caption, para que cada documento lo envuelva a su manera (`table` en
% la memoria, `frame` en beamer).
%
% Requiere: amsmath, booktabs.
% =====================================================================

\providecommand{\refcita}[1]{\texttt{#1}}

% ---------------------------------------------------------------------
% 1. Valores de los parámetros calibrados para Chile
% ---------------------------------------------------------------------
% Valores tal como los reporta el cuadro de parámetros base de la
% memoria. \valCAPdExacto es el valor sin redondear que reproducen las
% tablas de resultados; el capítulo de conclusiones usa 311,01.
% Pendiente de unificar por los autores: ver
% docs/auditoria/AUDITORIA_2_MHurtado.md.
\newcommand{\valCAPs}{264,20}          % MtCO2e
\newcommand{\valCAPd}{311,00}          % MtCO2e
\newcommand{\valCAPdExacto}{311,0134}  % MtCO2e
\newcommand{\valSigmaCAP}{61,86}       % MtCO2e
\newcommand{\valKappa}{64,40}          % USD/tCO2e
% Calibración: hay que distinguir dos valores de m que la memoria usa para
% cosas distintas (chapter04.tex, Cuadro tab:top10_distancia_mix y el texto
% que le sigue).
%   - El argmin de la distancia absoluta es m=0, que exige sigma_eps -> infinito
%     y por eso no sirve como parámetro operativo.
%   - Para los análisis posteriores se selecciona m=0,1, el menor m con un
%     sigma_eps finito, cuya distancia (0,2758) queda muy cerca del mínimo.
% Nota: el texto de la memoria dice sigma_eps = 186,90 y el cuadro dice 185,58
% para el mismo m=0,1. Pendiente de unificar por los autores; ver
% docs/auditoria/AUDITORIA_2_MHurtado.md.
\newcommand{\valMArgmin}{0}                  % menor distancia absoluta
\newcommand{\valDistanciaArgmin}{0,2727}
\newcommand{\valMSeleccionado}{0,1}          % usado en los análisis posteriores
\newcommand{\valDistanciaSeleccionado}{0,2758}
\newcommand{\valSigmaEpsSeleccionado}{185,58} % MtCO2e, segun tab:top10_distancia_mix
\newcommand{\valHorizonte}{2020--2030}
\newcommand{\unidadMt}{MtCO$_2$e}
\newcommand{\unidadUSDt}{USD/tCO$_2$e}

% ---------------------------------------------------------------------
% 2. Ecuaciones del planificador con inatención conductual
% ---------------------------------------------------------------------
% Cuerpo matemático sin entorno, para poder usarlo en `equation`,
% en \[...\] o dentro de un nodo tikz.
\newcommand{\ecSenal}{CAP_s = CAP + \epsilon}
\newcommand{\ecEsperanza}{\mathbb{E}\left[CAP|CAP_s\right] = m CAP_s + (1 - m)CAP_d}
\newcommand{\ecAtencion}{m = \frac{\sigma_{CAP}^2}{\sigma_{CAP}^2 + \sigma_\varepsilon^2}}
\newcommand{\ecThetaSinPrecio}{{\theta^*}(CAP_s) = m\cdot CAP_s + (1 - m)\cdot CAP_d}
\newcommand{\ecUtilidadPlanificador}{\mathcal{V}(\theta,CAP) = \pi^a \theta - \frac{\kappa}{2}(\theta - CAP)^2}
\newcommand{\ecThetaOptimo}{\theta^*(CAP_s) = m\,CAP_s + (1-m)\,CAP_d + \frac{\pi^a}{\kappa}}

% ---------------------------------------------------------------------
% 3. Descripciones de los símbolos, definidas una sola vez
% ---------------------------------------------------------------------
% La memoria las usa como filas de tabla y la presentación como
% tarjetas; por eso el texto se define aquí y el formato en cada
% documento.
\newcommand{\descCAP}{Presupuesto ambiental real}
\newcommand{\descCAPs}{Señal observada del presupuesto de emisiones}
\newcommand{\descCAPd}{Presupuesto determinista del CAP (valor por defecto)}
\newcommand{\descEps}{Ruido en la señal}
\newcommand{\descSigmaEps}{Varianza del error en la señal}
\newcommand{\descSigmaCAP}{Varianza del CAP}
\newcommand{\descTheta}{Permisos totales emitidos por el planificador}
\newcommand{\descPiA}{Precio de permisos emitidos por el planificador}
\newcommand{\descPiV}{Precio de permisos en el mercado entre productores}
\newcommand{\descPiD}{Precio de la electricidad pagado por el consumidor}
\newcommand{\descKappa}{Costo social por alejarse del $CAP$}
\newcommand{\descM}{Nivel de atención del planificador}
% Distribuciones
\newcommand{\distCAP}{CAP \sim \mathcal{N}(CAP_d,\sigma^2_{CAP})}
\newcommand{\distEps}{\epsilon \sim \mathcal{N}(0,\sigma^2_\varepsilon)}

% ---------------------------------------------------------------------
% 4. Nomenclatura del modelo de Amigo et al. (2021)
% ---------------------------------------------------------------------
\newcommand{\tablaNomenclaturaAmigo}{%
\begin{tabular}{lll}
\toprule
\textbf{Símbolo} & \textbf{Descripción} & \textbf{Unidad} \\
\midrule
$\pi^{a}$ & \descPiA & USD/$tCO_2e$ \\
$\pi^{v}$ & \descPiV & USD/permiso \\
$\pi^{d}$ & \descPiD & USD/MWh \\
\midrule
$x_i$ & Inversión o expansión de capacidad del productor $i$ & MW \\
$Q_i$ & Generación eléctrica del productor $i$ & MWh \\
$A_i$ & Permisos adquiridos por el productor $i$ al planificador & $tCO_2e$\\
$P_i$ & Permisos comprados por el productor $i$ en el mercado secundario & $tCO_2e$ \\
$V_i$ & Permisos vendidos por el productor $i$ en el mercado secundario & $tCO_2e$ \\
$\epsilon_i$ & Factor de emisión por tecnología $i$ & $tCO_2$e/MWh \\
\midrule
$\theta$ & \descTheta & $tCO_2e$ \\
\bottomrule
\end{tabular}}

% ---------------------------------------------------------------------
% 5. Parámetros que el modelo conductual agrega a Amigo et al. (2021)
% ---------------------------------------------------------------------
\newcommand{\tablaParametrosAgregados}{%
\begin{tabular}{lll}
\toprule
\textbf{Símbolo} & \textbf{Descripción} & \textbf{Unidad} \\
\midrule
$\kappa$                & \descKappa      & \unidadUSDt \\
$CAP_s$                 & \descCAPs       & tCO$_2$e \\
$CAP_d$                 & \descCAPd       & tCO$_2$e \\
$\sigma_\varepsilon^2$  & \descSigmaEps   & tCO$_2$e \\
$\sigma_{CAP}^2$        & \descSigmaCAP   & tCO$_2$e \\
\bottomrule
\end{tabular}}

% ---------------------------------------------------------------------
% 6. Valores base de los parámetros, con su origen
% ---------------------------------------------------------------------
% Las descripciones de origen se definen aparte para que la
% presentación pueda usarlas en tarjetas.
\newcommand{\origenCAPs}{\refcita{ministerio_de_energia_de_chile_plan_2024}}
\newcommand{\origenCAPd}{Promedio sector energía 2012--2022 (SNICHILE)}
\newcommand{\origenSigmaCAP}{Incertidumbre estructural histórica}
\newcommand{\origenKappa}{Costo social del carbono, Gobierno de Chile (2024)}

\newcommand{\tablaParametrosChile}{%
\begin{tabular}{l c l l}
\toprule
\textbf{Parámetro} & \textbf{Valor} & \textbf{Unidad} & \textbf{Origen} \\
\midrule
$CAP_s$          & \valCAPs      & \unidadMt   & \origenCAPs \\
$CAP_d$          & \valCAPd      & \unidadMt   & \origenCAPd \\
$\sigma_{CAP}$   & \valSigmaCAP  & \unidadMt   & \origenSigmaCAP \\
$\kappa$         & \valKappa     & \unidadUSDt & \origenKappa \\
\bottomrule
\end{tabular}}

% o
%   \newcommand{\refcita}[1]{\citet{#1}}   % presentación (natbib)
%   % =====================================================================
% shared/datos-modelo.tex
% ---------------------------------------------------------------------
% Nomenclatura, ecuaciones y valores de parámetros del modelo de
% anclaje y ajuste del subastador. Fuente única para la memoria
% (docs/MemoriaMHurtado) y la presentación de defensa
% (docs/Presentations/presentation-template), de modo que un valor o una
% ecuación se escriban una sola vez.
%
% Uso:
%   \newcommand{\refcita}[1]{\citeA{#1}}   % memoria (newapa/apacite)
%   \input{../../shared/datos-modelo}
% o
%   \newcommand{\refcita}[1]{\citet{#1}}   % presentación (natbib)
%   \input{../../../shared/datos-modelo}
%
% Los macros de tabla entregan sólo el `tabular`, sin flotante ni
% caption, para que cada documento lo envuelva a su manera (`table` en
% la memoria, `frame` en beamer).
%
% Requiere: amsmath, booktabs.
% =====================================================================

\providecommand{\refcita}[1]{\texttt{#1}}

% ---------------------------------------------------------------------
% 1. Valores de los parámetros calibrados para Chile
% ---------------------------------------------------------------------
% Valores tal como los reporta el cuadro de parámetros base de la
% memoria. \valCAPdExacto es el valor sin redondear que reproducen las
% tablas de resultados; el capítulo de conclusiones usa 311,01.
% Pendiente de unificar por los autores: ver
% docs/auditoria/AUDITORIA_2_MHurtado.md.
\newcommand{\valCAPs}{264,20}          % MtCO2e
\newcommand{\valCAPd}{311,00}          % MtCO2e
\newcommand{\valCAPdExacto}{311,0134}  % MtCO2e
\newcommand{\valSigmaCAP}{61,86}       % MtCO2e
\newcommand{\valKappa}{64,40}          % USD/tCO2e
% Calibración: hay que distinguir dos valores de m que la memoria usa para
% cosas distintas (chapter04.tex, Cuadro tab:top10_distancia_mix y el texto
% que le sigue).
%   - El argmin de la distancia absoluta es m=0, que exige sigma_eps -> infinito
%     y por eso no sirve como parámetro operativo.
%   - Para los análisis posteriores se selecciona m=0,1, el menor m con un
%     sigma_eps finito, cuya distancia (0,2758) queda muy cerca del mínimo.
% Nota: el texto de la memoria dice sigma_eps = 186,90 y el cuadro dice 185,58
% para el mismo m=0,1. Pendiente de unificar por los autores; ver
% docs/auditoria/AUDITORIA_2_MHurtado.md.
\newcommand{\valMArgmin}{0}                  % menor distancia absoluta
\newcommand{\valDistanciaArgmin}{0,2727}
\newcommand{\valMSeleccionado}{0,1}          % usado en los análisis posteriores
\newcommand{\valDistanciaSeleccionado}{0,2758}
\newcommand{\valSigmaEpsSeleccionado}{185,58} % MtCO2e, segun tab:top10_distancia_mix
\newcommand{\valHorizonte}{2020--2030}
\newcommand{\unidadMt}{MtCO$_2$e}
\newcommand{\unidadUSDt}{USD/tCO$_2$e}

% ---------------------------------------------------------------------
% 2. Ecuaciones del planificador con inatención conductual
% ---------------------------------------------------------------------
% Cuerpo matemático sin entorno, para poder usarlo en `equation`,
% en \[...\] o dentro de un nodo tikz.
\newcommand{\ecSenal}{CAP_s = CAP + \epsilon}
\newcommand{\ecEsperanza}{\mathbb{E}\left[CAP|CAP_s\right] = m CAP_s + (1 - m)CAP_d}
\newcommand{\ecAtencion}{m = \frac{\sigma_{CAP}^2}{\sigma_{CAP}^2 + \sigma_\varepsilon^2}}
\newcommand{\ecThetaSinPrecio}{{\theta^*}(CAP_s) = m\cdot CAP_s + (1 - m)\cdot CAP_d}
\newcommand{\ecUtilidadPlanificador}{\mathcal{V}(\theta,CAP) = \pi^a \theta - \frac{\kappa}{2}(\theta - CAP)^2}
\newcommand{\ecThetaOptimo}{\theta^*(CAP_s) = m\,CAP_s + (1-m)\,CAP_d + \frac{\pi^a}{\kappa}}

% ---------------------------------------------------------------------
% 3. Descripciones de los símbolos, definidas una sola vez
% ---------------------------------------------------------------------
% La memoria las usa como filas de tabla y la presentación como
% tarjetas; por eso el texto se define aquí y el formato en cada
% documento.
\newcommand{\descCAP}{Presupuesto ambiental real}
\newcommand{\descCAPs}{Señal observada del presupuesto de emisiones}
\newcommand{\descCAPd}{Presupuesto determinista del CAP (valor por defecto)}
\newcommand{\descEps}{Ruido en la señal}
\newcommand{\descSigmaEps}{Varianza del error en la señal}
\newcommand{\descSigmaCAP}{Varianza del CAP}
\newcommand{\descTheta}{Permisos totales emitidos por el planificador}
\newcommand{\descPiA}{Precio de permisos emitidos por el planificador}
\newcommand{\descPiV}{Precio de permisos en el mercado entre productores}
\newcommand{\descPiD}{Precio de la electricidad pagado por el consumidor}
\newcommand{\descKappa}{Costo social por alejarse del $CAP$}
\newcommand{\descM}{Nivel de atención del planificador}
% Distribuciones
\newcommand{\distCAP}{CAP \sim \mathcal{N}(CAP_d,\sigma^2_{CAP})}
\newcommand{\distEps}{\epsilon \sim \mathcal{N}(0,\sigma^2_\varepsilon)}

% ---------------------------------------------------------------------
% 4. Nomenclatura del modelo de Amigo et al. (2021)
% ---------------------------------------------------------------------
\newcommand{\tablaNomenclaturaAmigo}{%
\begin{tabular}{lll}
\toprule
\textbf{Símbolo} & \textbf{Descripción} & \textbf{Unidad} \\
\midrule
$\pi^{a}$ & \descPiA & USD/$tCO_2e$ \\
$\pi^{v}$ & \descPiV & USD/permiso \\
$\pi^{d}$ & \descPiD & USD/MWh \\
\midrule
$x_i$ & Inversión o expansión de capacidad del productor $i$ & MW \\
$Q_i$ & Generación eléctrica del productor $i$ & MWh \\
$A_i$ & Permisos adquiridos por el productor $i$ al planificador & $tCO_2e$\\
$P_i$ & Permisos comprados por el productor $i$ en el mercado secundario & $tCO_2e$ \\
$V_i$ & Permisos vendidos por el productor $i$ en el mercado secundario & $tCO_2e$ \\
$\epsilon_i$ & Factor de emisión por tecnología $i$ & $tCO_2$e/MWh \\
\midrule
$\theta$ & \descTheta & $tCO_2e$ \\
\bottomrule
\end{tabular}}

% ---------------------------------------------------------------------
% 5. Parámetros que el modelo conductual agrega a Amigo et al. (2021)
% ---------------------------------------------------------------------
\newcommand{\tablaParametrosAgregados}{%
\begin{tabular}{lll}
\toprule
\textbf{Símbolo} & \textbf{Descripción} & \textbf{Unidad} \\
\midrule
$\kappa$                & \descKappa      & \unidadUSDt \\
$CAP_s$                 & \descCAPs       & tCO$_2$e \\
$CAP_d$                 & \descCAPd       & tCO$_2$e \\
$\sigma_\varepsilon^2$  & \descSigmaEps   & tCO$_2$e \\
$\sigma_{CAP}^2$        & \descSigmaCAP   & tCO$_2$e \\
\bottomrule
\end{tabular}}

% ---------------------------------------------------------------------
% 6. Valores base de los parámetros, con su origen
% ---------------------------------------------------------------------
% Las descripciones de origen se definen aparte para que la
% presentación pueda usarlas en tarjetas.
\newcommand{\origenCAPs}{\refcita{ministerio_de_energia_de_chile_plan_2024}}
\newcommand{\origenCAPd}{Promedio sector energía 2012--2022 (SNICHILE)}
\newcommand{\origenSigmaCAP}{Incertidumbre estructural histórica}
\newcommand{\origenKappa}{Costo social del carbono, Gobierno de Chile (2024)}

\newcommand{\tablaParametrosChile}{%
\begin{tabular}{l c l l}
\toprule
\textbf{Parámetro} & \textbf{Valor} & \textbf{Unidad} & \textbf{Origen} \\
\midrule
$CAP_s$          & \valCAPs      & \unidadMt   & \origenCAPs \\
$CAP_d$          & \valCAPd      & \unidadMt   & \origenCAPd \\
$\sigma_{CAP}$   & \valSigmaCAP  & \unidadMt   & \origenSigmaCAP \\
$\kappa$         & \valKappa     & \unidadUSDt & \origenKappa \\
\bottomrule
\end{tabular}}

%
% Los macros de tabla entregan sólo el `tabular`, sin flotante ni
% caption, para que cada documento lo envuelva a su manera (`table` en
% la memoria, `frame` en beamer).
%
% Requiere: amsmath, booktabs.
% =====================================================================

\providecommand{\refcita}[1]{\texttt{#1}}

% ---------------------------------------------------------------------
% 1. Valores de los parámetros calibrados para Chile
% ---------------------------------------------------------------------
% Valores tal como los reporta el cuadro de parámetros base de la
% memoria. \valCAPdExacto es el valor sin redondear que reproducen las
% tablas de resultados; el capítulo de conclusiones usa 311,01.
% Pendiente de unificar por los autores: ver
% docs/auditoria/AUDITORIA_2_MHurtado.md.
\newcommand{\valCAPs}{264,20}          % MtCO2e
\newcommand{\valCAPd}{311,00}          % MtCO2e
\newcommand{\valCAPdExacto}{311,0134}  % MtCO2e
\newcommand{\valSigmaCAP}{61,86}       % MtCO2e
\newcommand{\valKappa}{64,40}          % USD/tCO2e
% Calibración: hay que distinguir dos valores de m que la memoria usa para
% cosas distintas (chapter04.tex, Cuadro tab:top10_distancia_mix y el texto
% que le sigue).
%   - El argmin de la distancia absoluta es m=0, que exige sigma_eps -> infinito
%     y por eso no sirve como parámetro operativo.
%   - Para los análisis posteriores se selecciona m=0,1, el menor m con un
%     sigma_eps finito, cuya distancia (0,2758) queda muy cerca del mínimo.
% Nota: el texto de la memoria dice sigma_eps = 186,90 y el cuadro dice 185,58
% para el mismo m=0,1. Pendiente de unificar por los autores; ver
% docs/auditoria/AUDITORIA_2_MHurtado.md.
\newcommand{\valMArgmin}{0}                  % menor distancia absoluta
\newcommand{\valDistanciaArgmin}{0,2727}
\newcommand{\valMSeleccionado}{0,1}          % usado en los análisis posteriores
\newcommand{\valDistanciaSeleccionado}{0,2758}
\newcommand{\valSigmaEpsSeleccionado}{185,58} % MtCO2e, segun tab:top10_distancia_mix
\newcommand{\valHorizonte}{2020--2030}
\newcommand{\unidadMt}{MtCO$_2$e}
\newcommand{\unidadUSDt}{USD/tCO$_2$e}

% ---------------------------------------------------------------------
% 2. Ecuaciones del planificador con inatención conductual
% ---------------------------------------------------------------------
% Cuerpo matemático sin entorno, para poder usarlo en `equation`,
% en \[...\] o dentro de un nodo tikz.
\newcommand{\ecSenal}{CAP_s = CAP + \epsilon}
\newcommand{\ecEsperanza}{\mathbb{E}\left[CAP|CAP_s\right] = m CAP_s + (1 - m)CAP_d}
\newcommand{\ecAtencion}{m = \frac{\sigma_{CAP}^2}{\sigma_{CAP}^2 + \sigma_\varepsilon^2}}
\newcommand{\ecThetaSinPrecio}{{\theta^*}(CAP_s) = m\cdot CAP_s + (1 - m)\cdot CAP_d}
\newcommand{\ecUtilidadPlanificador}{\mathcal{V}(\theta,CAP) = \pi^a \theta - \frac{\kappa}{2}(\theta - CAP)^2}
\newcommand{\ecThetaOptimo}{\theta^*(CAP_s) = m\,CAP_s + (1-m)\,CAP_d + \frac{\pi^a}{\kappa}}

% ---------------------------------------------------------------------
% 3. Descripciones de los símbolos, definidas una sola vez
% ---------------------------------------------------------------------
% La memoria las usa como filas de tabla y la presentación como
% tarjetas; por eso el texto se define aquí y el formato en cada
% documento.
\newcommand{\descCAP}{Presupuesto ambiental real}
\newcommand{\descCAPs}{Señal observada del presupuesto de emisiones}
\newcommand{\descCAPd}{Presupuesto determinista del CAP (valor por defecto)}
\newcommand{\descEps}{Ruido en la señal}
\newcommand{\descSigmaEps}{Varianza del error en la señal}
\newcommand{\descSigmaCAP}{Varianza del CAP}
\newcommand{\descTheta}{Permisos totales emitidos por el planificador}
\newcommand{\descPiA}{Precio de permisos emitidos por el planificador}
\newcommand{\descPiV}{Precio de permisos en el mercado entre productores}
\newcommand{\descPiD}{Precio de la electricidad pagado por el consumidor}
\newcommand{\descKappa}{Costo social por alejarse del $CAP$}
\newcommand{\descM}{Nivel de atención del planificador}
% Distribuciones
\newcommand{\distCAP}{CAP \sim \mathcal{N}(CAP_d,\sigma^2_{CAP})}
\newcommand{\distEps}{\epsilon \sim \mathcal{N}(0,\sigma^2_\varepsilon)}

% ---------------------------------------------------------------------
% 4. Nomenclatura del modelo de Amigo et al. (2021)
% ---------------------------------------------------------------------
\newcommand{\tablaNomenclaturaAmigo}{%
\begin{tabular}{lll}
\toprule
\textbf{Símbolo} & \textbf{Descripción} & \textbf{Unidad} \\
\midrule
$\pi^{a}$ & \descPiA & USD/$tCO_2e$ \\
$\pi^{v}$ & \descPiV & USD/permiso \\
$\pi^{d}$ & \descPiD & USD/MWh \\
\midrule
$x_i$ & Inversión o expansión de capacidad del productor $i$ & MW \\
$Q_i$ & Generación eléctrica del productor $i$ & MWh \\
$A_i$ & Permisos adquiridos por el productor $i$ al planificador & $tCO_2e$\\
$P_i$ & Permisos comprados por el productor $i$ en el mercado secundario & $tCO_2e$ \\
$V_i$ & Permisos vendidos por el productor $i$ en el mercado secundario & $tCO_2e$ \\
$\epsilon_i$ & Factor de emisión por tecnología $i$ & $tCO_2$e/MWh \\
\midrule
$\theta$ & \descTheta & $tCO_2e$ \\
\bottomrule
\end{tabular}}

% ---------------------------------------------------------------------
% 5. Parámetros que el modelo conductual agrega a Amigo et al. (2021)
% ---------------------------------------------------------------------
\newcommand{\tablaParametrosAgregados}{%
\begin{tabular}{lll}
\toprule
\textbf{Símbolo} & \textbf{Descripción} & \textbf{Unidad} \\
\midrule
$\kappa$                & \descKappa      & \unidadUSDt \\
$CAP_s$                 & \descCAPs       & tCO$_2$e \\
$CAP_d$                 & \descCAPd       & tCO$_2$e \\
$\sigma_\varepsilon^2$  & \descSigmaEps   & tCO$_2$e \\
$\sigma_{CAP}^2$        & \descSigmaCAP   & tCO$_2$e \\
\bottomrule
\end{tabular}}

% ---------------------------------------------------------------------
% 6. Valores base de los parámetros, con su origen
% ---------------------------------------------------------------------
% Las descripciones de origen se definen aparte para que la
% presentación pueda usarlas en tarjetas.
\newcommand{\origenCAPs}{\refcita{ministerio_de_energia_de_chile_plan_2024}}
\newcommand{\origenCAPd}{Promedio sector energía 2012--2022 (SNICHILE)}
\newcommand{\origenSigmaCAP}{Incertidumbre estructural histórica}
\newcommand{\origenKappa}{Costo social del carbono, Gobierno de Chile (2024)}

\newcommand{\tablaParametrosChile}{%
\begin{tabular}{l c l l}
\toprule
\textbf{Parámetro} & \textbf{Valor} & \textbf{Unidad} & \textbf{Origen} \\
\midrule
$CAP_s$          & \valCAPs      & \unidadMt   & \origenCAPs \\
$CAP_d$          & \valCAPd      & \unidadMt   & \origenCAPd \\
$\sigma_{CAP}$   & \valSigmaCAP  & \unidadMt   & \origenSigmaCAP \\
$\kappa$         & \valKappa     & \unidadUSDt & \origenKappa \\
\bottomrule
\end{tabular}}

% o
%   \newcommand{\refcita}[1]{\citet{#1}}   % presentación (natbib)
%   % =====================================================================
% shared/datos-modelo.tex
% ---------------------------------------------------------------------
% Nomenclatura, ecuaciones y valores de parámetros del modelo de
% anclaje y ajuste del subastador. Fuente única para la memoria
% (docs/MemoriaMHurtado) y la presentación de defensa
% (docs/Presentations/presentation-template), de modo que un valor o una
% ecuación se escriban una sola vez.
%
% Uso:
%   \newcommand{\refcita}[1]{\citeA{#1}}   % memoria (newapa/apacite)
%   % =====================================================================
% shared/datos-modelo.tex
% ---------------------------------------------------------------------
% Nomenclatura, ecuaciones y valores de parámetros del modelo de
% anclaje y ajuste del subastador. Fuente única para la memoria
% (docs/MemoriaMHurtado) y la presentación de defensa
% (docs/Presentations/presentation-template), de modo que un valor o una
% ecuación se escriban una sola vez.
%
% Uso:
%   \newcommand{\refcita}[1]{\citeA{#1}}   % memoria (newapa/apacite)
%   \input{../../shared/datos-modelo}
% o
%   \newcommand{\refcita}[1]{\citet{#1}}   % presentación (natbib)
%   \input{../../../shared/datos-modelo}
%
% Los macros de tabla entregan sólo el `tabular`, sin flotante ni
% caption, para que cada documento lo envuelva a su manera (`table` en
% la memoria, `frame` en beamer).
%
% Requiere: amsmath, booktabs.
% =====================================================================

\providecommand{\refcita}[1]{\texttt{#1}}

% ---------------------------------------------------------------------
% 1. Valores de los parámetros calibrados para Chile
% ---------------------------------------------------------------------
% Valores tal como los reporta el cuadro de parámetros base de la
% memoria. \valCAPdExacto es el valor sin redondear que reproducen las
% tablas de resultados; el capítulo de conclusiones usa 311,01.
% Pendiente de unificar por los autores: ver
% docs/auditoria/AUDITORIA_2_MHurtado.md.
\newcommand{\valCAPs}{264,20}          % MtCO2e
\newcommand{\valCAPd}{311,00}          % MtCO2e
\newcommand{\valCAPdExacto}{311,0134}  % MtCO2e
\newcommand{\valSigmaCAP}{61,86}       % MtCO2e
\newcommand{\valKappa}{64,40}          % USD/tCO2e
% Calibración: hay que distinguir dos valores de m que la memoria usa para
% cosas distintas (chapter04.tex, Cuadro tab:top10_distancia_mix y el texto
% que le sigue).
%   - El argmin de la distancia absoluta es m=0, que exige sigma_eps -> infinito
%     y por eso no sirve como parámetro operativo.
%   - Para los análisis posteriores se selecciona m=0,1, el menor m con un
%     sigma_eps finito, cuya distancia (0,2758) queda muy cerca del mínimo.
% Nota: el texto de la memoria dice sigma_eps = 186,90 y el cuadro dice 185,58
% para el mismo m=0,1. Pendiente de unificar por los autores; ver
% docs/auditoria/AUDITORIA_2_MHurtado.md.
\newcommand{\valMArgmin}{0}                  % menor distancia absoluta
\newcommand{\valDistanciaArgmin}{0,2727}
\newcommand{\valMSeleccionado}{0,1}          % usado en los análisis posteriores
\newcommand{\valDistanciaSeleccionado}{0,2758}
\newcommand{\valSigmaEpsSeleccionado}{185,58} % MtCO2e, segun tab:top10_distancia_mix
\newcommand{\valHorizonte}{2020--2030}
\newcommand{\unidadMt}{MtCO$_2$e}
\newcommand{\unidadUSDt}{USD/tCO$_2$e}

% ---------------------------------------------------------------------
% 2. Ecuaciones del planificador con inatención conductual
% ---------------------------------------------------------------------
% Cuerpo matemático sin entorno, para poder usarlo en `equation`,
% en \[...\] o dentro de un nodo tikz.
\newcommand{\ecSenal}{CAP_s = CAP + \epsilon}
\newcommand{\ecEsperanza}{\mathbb{E}\left[CAP|CAP_s\right] = m CAP_s + (1 - m)CAP_d}
\newcommand{\ecAtencion}{m = \frac{\sigma_{CAP}^2}{\sigma_{CAP}^2 + \sigma_\varepsilon^2}}
\newcommand{\ecThetaSinPrecio}{{\theta^*}(CAP_s) = m\cdot CAP_s + (1 - m)\cdot CAP_d}
\newcommand{\ecUtilidadPlanificador}{\mathcal{V}(\theta,CAP) = \pi^a \theta - \frac{\kappa}{2}(\theta - CAP)^2}
\newcommand{\ecThetaOptimo}{\theta^*(CAP_s) = m\,CAP_s + (1-m)\,CAP_d + \frac{\pi^a}{\kappa}}

% ---------------------------------------------------------------------
% 3. Descripciones de los símbolos, definidas una sola vez
% ---------------------------------------------------------------------
% La memoria las usa como filas de tabla y la presentación como
% tarjetas; por eso el texto se define aquí y el formato en cada
% documento.
\newcommand{\descCAP}{Presupuesto ambiental real}
\newcommand{\descCAPs}{Señal observada del presupuesto de emisiones}
\newcommand{\descCAPd}{Presupuesto determinista del CAP (valor por defecto)}
\newcommand{\descEps}{Ruido en la señal}
\newcommand{\descSigmaEps}{Varianza del error en la señal}
\newcommand{\descSigmaCAP}{Varianza del CAP}
\newcommand{\descTheta}{Permisos totales emitidos por el planificador}
\newcommand{\descPiA}{Precio de permisos emitidos por el planificador}
\newcommand{\descPiV}{Precio de permisos en el mercado entre productores}
\newcommand{\descPiD}{Precio de la electricidad pagado por el consumidor}
\newcommand{\descKappa}{Costo social por alejarse del $CAP$}
\newcommand{\descM}{Nivel de atención del planificador}
% Distribuciones
\newcommand{\distCAP}{CAP \sim \mathcal{N}(CAP_d,\sigma^2_{CAP})}
\newcommand{\distEps}{\epsilon \sim \mathcal{N}(0,\sigma^2_\varepsilon)}

% ---------------------------------------------------------------------
% 4. Nomenclatura del modelo de Amigo et al. (2021)
% ---------------------------------------------------------------------
\newcommand{\tablaNomenclaturaAmigo}{%
\begin{tabular}{lll}
\toprule
\textbf{Símbolo} & \textbf{Descripción} & \textbf{Unidad} \\
\midrule
$\pi^{a}$ & \descPiA & USD/$tCO_2e$ \\
$\pi^{v}$ & \descPiV & USD/permiso \\
$\pi^{d}$ & \descPiD & USD/MWh \\
\midrule
$x_i$ & Inversión o expansión de capacidad del productor $i$ & MW \\
$Q_i$ & Generación eléctrica del productor $i$ & MWh \\
$A_i$ & Permisos adquiridos por el productor $i$ al planificador & $tCO_2e$\\
$P_i$ & Permisos comprados por el productor $i$ en el mercado secundario & $tCO_2e$ \\
$V_i$ & Permisos vendidos por el productor $i$ en el mercado secundario & $tCO_2e$ \\
$\epsilon_i$ & Factor de emisión por tecnología $i$ & $tCO_2$e/MWh \\
\midrule
$\theta$ & \descTheta & $tCO_2e$ \\
\bottomrule
\end{tabular}}

% ---------------------------------------------------------------------
% 5. Parámetros que el modelo conductual agrega a Amigo et al. (2021)
% ---------------------------------------------------------------------
\newcommand{\tablaParametrosAgregados}{%
\begin{tabular}{lll}
\toprule
\textbf{Símbolo} & \textbf{Descripción} & \textbf{Unidad} \\
\midrule
$\kappa$                & \descKappa      & \unidadUSDt \\
$CAP_s$                 & \descCAPs       & tCO$_2$e \\
$CAP_d$                 & \descCAPd       & tCO$_2$e \\
$\sigma_\varepsilon^2$  & \descSigmaEps   & tCO$_2$e \\
$\sigma_{CAP}^2$        & \descSigmaCAP   & tCO$_2$e \\
\bottomrule
\end{tabular}}

% ---------------------------------------------------------------------
% 6. Valores base de los parámetros, con su origen
% ---------------------------------------------------------------------
% Las descripciones de origen se definen aparte para que la
% presentación pueda usarlas en tarjetas.
\newcommand{\origenCAPs}{\refcita{ministerio_de_energia_de_chile_plan_2024}}
\newcommand{\origenCAPd}{Promedio sector energía 2012--2022 (SNICHILE)}
\newcommand{\origenSigmaCAP}{Incertidumbre estructural histórica}
\newcommand{\origenKappa}{Costo social del carbono, Gobierno de Chile (2024)}

\newcommand{\tablaParametrosChile}{%
\begin{tabular}{l c l l}
\toprule
\textbf{Parámetro} & \textbf{Valor} & \textbf{Unidad} & \textbf{Origen} \\
\midrule
$CAP_s$          & \valCAPs      & \unidadMt   & \origenCAPs \\
$CAP_d$          & \valCAPd      & \unidadMt   & \origenCAPd \\
$\sigma_{CAP}$   & \valSigmaCAP  & \unidadMt   & \origenSigmaCAP \\
$\kappa$         & \valKappa     & \unidadUSDt & \origenKappa \\
\bottomrule
\end{tabular}}

% o
%   \newcommand{\refcita}[1]{\citet{#1}}   % presentación (natbib)
%   % =====================================================================
% shared/datos-modelo.tex
% ---------------------------------------------------------------------
% Nomenclatura, ecuaciones y valores de parámetros del modelo de
% anclaje y ajuste del subastador. Fuente única para la memoria
% (docs/MemoriaMHurtado) y la presentación de defensa
% (docs/Presentations/presentation-template), de modo que un valor o una
% ecuación se escriban una sola vez.
%
% Uso:
%   \newcommand{\refcita}[1]{\citeA{#1}}   % memoria (newapa/apacite)
%   \input{../../shared/datos-modelo}
% o
%   \newcommand{\refcita}[1]{\citet{#1}}   % presentación (natbib)
%   \input{../../../shared/datos-modelo}
%
% Los macros de tabla entregan sólo el `tabular`, sin flotante ni
% caption, para que cada documento lo envuelva a su manera (`table` en
% la memoria, `frame` en beamer).
%
% Requiere: amsmath, booktabs.
% =====================================================================

\providecommand{\refcita}[1]{\texttt{#1}}

% ---------------------------------------------------------------------
% 1. Valores de los parámetros calibrados para Chile
% ---------------------------------------------------------------------
% Valores tal como los reporta el cuadro de parámetros base de la
% memoria. \valCAPdExacto es el valor sin redondear que reproducen las
% tablas de resultados; el capítulo de conclusiones usa 311,01.
% Pendiente de unificar por los autores: ver
% docs/auditoria/AUDITORIA_2_MHurtado.md.
\newcommand{\valCAPs}{264,20}          % MtCO2e
\newcommand{\valCAPd}{311,00}          % MtCO2e
\newcommand{\valCAPdExacto}{311,0134}  % MtCO2e
\newcommand{\valSigmaCAP}{61,86}       % MtCO2e
\newcommand{\valKappa}{64,40}          % USD/tCO2e
% Calibración: hay que distinguir dos valores de m que la memoria usa para
% cosas distintas (chapter04.tex, Cuadro tab:top10_distancia_mix y el texto
% que le sigue).
%   - El argmin de la distancia absoluta es m=0, que exige sigma_eps -> infinito
%     y por eso no sirve como parámetro operativo.
%   - Para los análisis posteriores se selecciona m=0,1, el menor m con un
%     sigma_eps finito, cuya distancia (0,2758) queda muy cerca del mínimo.
% Nota: el texto de la memoria dice sigma_eps = 186,90 y el cuadro dice 185,58
% para el mismo m=0,1. Pendiente de unificar por los autores; ver
% docs/auditoria/AUDITORIA_2_MHurtado.md.
\newcommand{\valMArgmin}{0}                  % menor distancia absoluta
\newcommand{\valDistanciaArgmin}{0,2727}
\newcommand{\valMSeleccionado}{0,1}          % usado en los análisis posteriores
\newcommand{\valDistanciaSeleccionado}{0,2758}
\newcommand{\valSigmaEpsSeleccionado}{185,58} % MtCO2e, segun tab:top10_distancia_mix
\newcommand{\valHorizonte}{2020--2030}
\newcommand{\unidadMt}{MtCO$_2$e}
\newcommand{\unidadUSDt}{USD/tCO$_2$e}

% ---------------------------------------------------------------------
% 2. Ecuaciones del planificador con inatención conductual
% ---------------------------------------------------------------------
% Cuerpo matemático sin entorno, para poder usarlo en `equation`,
% en \[...\] o dentro de un nodo tikz.
\newcommand{\ecSenal}{CAP_s = CAP + \epsilon}
\newcommand{\ecEsperanza}{\mathbb{E}\left[CAP|CAP_s\right] = m CAP_s + (1 - m)CAP_d}
\newcommand{\ecAtencion}{m = \frac{\sigma_{CAP}^2}{\sigma_{CAP}^2 + \sigma_\varepsilon^2}}
\newcommand{\ecThetaSinPrecio}{{\theta^*}(CAP_s) = m\cdot CAP_s + (1 - m)\cdot CAP_d}
\newcommand{\ecUtilidadPlanificador}{\mathcal{V}(\theta,CAP) = \pi^a \theta - \frac{\kappa}{2}(\theta - CAP)^2}
\newcommand{\ecThetaOptimo}{\theta^*(CAP_s) = m\,CAP_s + (1-m)\,CAP_d + \frac{\pi^a}{\kappa}}

% ---------------------------------------------------------------------
% 3. Descripciones de los símbolos, definidas una sola vez
% ---------------------------------------------------------------------
% La memoria las usa como filas de tabla y la presentación como
% tarjetas; por eso el texto se define aquí y el formato en cada
% documento.
\newcommand{\descCAP}{Presupuesto ambiental real}
\newcommand{\descCAPs}{Señal observada del presupuesto de emisiones}
\newcommand{\descCAPd}{Presupuesto determinista del CAP (valor por defecto)}
\newcommand{\descEps}{Ruido en la señal}
\newcommand{\descSigmaEps}{Varianza del error en la señal}
\newcommand{\descSigmaCAP}{Varianza del CAP}
\newcommand{\descTheta}{Permisos totales emitidos por el planificador}
\newcommand{\descPiA}{Precio de permisos emitidos por el planificador}
\newcommand{\descPiV}{Precio de permisos en el mercado entre productores}
\newcommand{\descPiD}{Precio de la electricidad pagado por el consumidor}
\newcommand{\descKappa}{Costo social por alejarse del $CAP$}
\newcommand{\descM}{Nivel de atención del planificador}
% Distribuciones
\newcommand{\distCAP}{CAP \sim \mathcal{N}(CAP_d,\sigma^2_{CAP})}
\newcommand{\distEps}{\epsilon \sim \mathcal{N}(0,\sigma^2_\varepsilon)}

% ---------------------------------------------------------------------
% 4. Nomenclatura del modelo de Amigo et al. (2021)
% ---------------------------------------------------------------------
\newcommand{\tablaNomenclaturaAmigo}{%
\begin{tabular}{lll}
\toprule
\textbf{Símbolo} & \textbf{Descripción} & \textbf{Unidad} \\
\midrule
$\pi^{a}$ & \descPiA & USD/$tCO_2e$ \\
$\pi^{v}$ & \descPiV & USD/permiso \\
$\pi^{d}$ & \descPiD & USD/MWh \\
\midrule
$x_i$ & Inversión o expansión de capacidad del productor $i$ & MW \\
$Q_i$ & Generación eléctrica del productor $i$ & MWh \\
$A_i$ & Permisos adquiridos por el productor $i$ al planificador & $tCO_2e$\\
$P_i$ & Permisos comprados por el productor $i$ en el mercado secundario & $tCO_2e$ \\
$V_i$ & Permisos vendidos por el productor $i$ en el mercado secundario & $tCO_2e$ \\
$\epsilon_i$ & Factor de emisión por tecnología $i$ & $tCO_2$e/MWh \\
\midrule
$\theta$ & \descTheta & $tCO_2e$ \\
\bottomrule
\end{tabular}}

% ---------------------------------------------------------------------
% 5. Parámetros que el modelo conductual agrega a Amigo et al. (2021)
% ---------------------------------------------------------------------
\newcommand{\tablaParametrosAgregados}{%
\begin{tabular}{lll}
\toprule
\textbf{Símbolo} & \textbf{Descripción} & \textbf{Unidad} \\
\midrule
$\kappa$                & \descKappa      & \unidadUSDt \\
$CAP_s$                 & \descCAPs       & tCO$_2$e \\
$CAP_d$                 & \descCAPd       & tCO$_2$e \\
$\sigma_\varepsilon^2$  & \descSigmaEps   & tCO$_2$e \\
$\sigma_{CAP}^2$        & \descSigmaCAP   & tCO$_2$e \\
\bottomrule
\end{tabular}}

% ---------------------------------------------------------------------
% 6. Valores base de los parámetros, con su origen
% ---------------------------------------------------------------------
% Las descripciones de origen se definen aparte para que la
% presentación pueda usarlas en tarjetas.
\newcommand{\origenCAPs}{\refcita{ministerio_de_energia_de_chile_plan_2024}}
\newcommand{\origenCAPd}{Promedio sector energía 2012--2022 (SNICHILE)}
\newcommand{\origenSigmaCAP}{Incertidumbre estructural histórica}
\newcommand{\origenKappa}{Costo social del carbono, Gobierno de Chile (2024)}

\newcommand{\tablaParametrosChile}{%
\begin{tabular}{l c l l}
\toprule
\textbf{Parámetro} & \textbf{Valor} & \textbf{Unidad} & \textbf{Origen} \\
\midrule
$CAP_s$          & \valCAPs      & \unidadMt   & \origenCAPs \\
$CAP_d$          & \valCAPd      & \unidadMt   & \origenCAPd \\
$\sigma_{CAP}$   & \valSigmaCAP  & \unidadMt   & \origenSigmaCAP \\
$\kappa$         & \valKappa     & \unidadUSDt & \origenKappa \\
\bottomrule
\end{tabular}}

%
% Los macros de tabla entregan sólo el `tabular`, sin flotante ni
% caption, para que cada documento lo envuelva a su manera (`table` en
% la memoria, `frame` en beamer).
%
% Requiere: amsmath, booktabs.
% =====================================================================

\providecommand{\refcita}[1]{\texttt{#1}}

% ---------------------------------------------------------------------
% 1. Valores de los parámetros calibrados para Chile
% ---------------------------------------------------------------------
% Valores tal como los reporta el cuadro de parámetros base de la
% memoria. \valCAPdExacto es el valor sin redondear que reproducen las
% tablas de resultados; el capítulo de conclusiones usa 311,01.
% Pendiente de unificar por los autores: ver
% docs/auditoria/AUDITORIA_2_MHurtado.md.
\newcommand{\valCAPs}{264,20}          % MtCO2e
\newcommand{\valCAPd}{311,00}          % MtCO2e
\newcommand{\valCAPdExacto}{311,0134}  % MtCO2e
\newcommand{\valSigmaCAP}{61,86}       % MtCO2e
\newcommand{\valKappa}{64,40}          % USD/tCO2e
% Calibración: hay que distinguir dos valores de m que la memoria usa para
% cosas distintas (chapter04.tex, Cuadro tab:top10_distancia_mix y el texto
% que le sigue).
%   - El argmin de la distancia absoluta es m=0, que exige sigma_eps -> infinito
%     y por eso no sirve como parámetro operativo.
%   - Para los análisis posteriores se selecciona m=0,1, el menor m con un
%     sigma_eps finito, cuya distancia (0,2758) queda muy cerca del mínimo.
% Nota: el texto de la memoria dice sigma_eps = 186,90 y el cuadro dice 185,58
% para el mismo m=0,1. Pendiente de unificar por los autores; ver
% docs/auditoria/AUDITORIA_2_MHurtado.md.
\newcommand{\valMArgmin}{0}                  % menor distancia absoluta
\newcommand{\valDistanciaArgmin}{0,2727}
\newcommand{\valMSeleccionado}{0,1}          % usado en los análisis posteriores
\newcommand{\valDistanciaSeleccionado}{0,2758}
\newcommand{\valSigmaEpsSeleccionado}{185,58} % MtCO2e, segun tab:top10_distancia_mix
\newcommand{\valHorizonte}{2020--2030}
\newcommand{\unidadMt}{MtCO$_2$e}
\newcommand{\unidadUSDt}{USD/tCO$_2$e}

% ---------------------------------------------------------------------
% 2. Ecuaciones del planificador con inatención conductual
% ---------------------------------------------------------------------
% Cuerpo matemático sin entorno, para poder usarlo en `equation`,
% en \[...\] o dentro de un nodo tikz.
\newcommand{\ecSenal}{CAP_s = CAP + \epsilon}
\newcommand{\ecEsperanza}{\mathbb{E}\left[CAP|CAP_s\right] = m CAP_s + (1 - m)CAP_d}
\newcommand{\ecAtencion}{m = \frac{\sigma_{CAP}^2}{\sigma_{CAP}^2 + \sigma_\varepsilon^2}}
\newcommand{\ecThetaSinPrecio}{{\theta^*}(CAP_s) = m\cdot CAP_s + (1 - m)\cdot CAP_d}
\newcommand{\ecUtilidadPlanificador}{\mathcal{V}(\theta,CAP) = \pi^a \theta - \frac{\kappa}{2}(\theta - CAP)^2}
\newcommand{\ecThetaOptimo}{\theta^*(CAP_s) = m\,CAP_s + (1-m)\,CAP_d + \frac{\pi^a}{\kappa}}

% ---------------------------------------------------------------------
% 3. Descripciones de los símbolos, definidas una sola vez
% ---------------------------------------------------------------------
% La memoria las usa como filas de tabla y la presentación como
% tarjetas; por eso el texto se define aquí y el formato en cada
% documento.
\newcommand{\descCAP}{Presupuesto ambiental real}
\newcommand{\descCAPs}{Señal observada del presupuesto de emisiones}
\newcommand{\descCAPd}{Presupuesto determinista del CAP (valor por defecto)}
\newcommand{\descEps}{Ruido en la señal}
\newcommand{\descSigmaEps}{Varianza del error en la señal}
\newcommand{\descSigmaCAP}{Varianza del CAP}
\newcommand{\descTheta}{Permisos totales emitidos por el planificador}
\newcommand{\descPiA}{Precio de permisos emitidos por el planificador}
\newcommand{\descPiV}{Precio de permisos en el mercado entre productores}
\newcommand{\descPiD}{Precio de la electricidad pagado por el consumidor}
\newcommand{\descKappa}{Costo social por alejarse del $CAP$}
\newcommand{\descM}{Nivel de atención del planificador}
% Distribuciones
\newcommand{\distCAP}{CAP \sim \mathcal{N}(CAP_d,\sigma^2_{CAP})}
\newcommand{\distEps}{\epsilon \sim \mathcal{N}(0,\sigma^2_\varepsilon)}

% ---------------------------------------------------------------------
% 4. Nomenclatura del modelo de Amigo et al. (2021)
% ---------------------------------------------------------------------
\newcommand{\tablaNomenclaturaAmigo}{%
\begin{tabular}{lll}
\toprule
\textbf{Símbolo} & \textbf{Descripción} & \textbf{Unidad} \\
\midrule
$\pi^{a}$ & \descPiA & USD/$tCO_2e$ \\
$\pi^{v}$ & \descPiV & USD/permiso \\
$\pi^{d}$ & \descPiD & USD/MWh \\
\midrule
$x_i$ & Inversión o expansión de capacidad del productor $i$ & MW \\
$Q_i$ & Generación eléctrica del productor $i$ & MWh \\
$A_i$ & Permisos adquiridos por el productor $i$ al planificador & $tCO_2e$\\
$P_i$ & Permisos comprados por el productor $i$ en el mercado secundario & $tCO_2e$ \\
$V_i$ & Permisos vendidos por el productor $i$ en el mercado secundario & $tCO_2e$ \\
$\epsilon_i$ & Factor de emisión por tecnología $i$ & $tCO_2$e/MWh \\
\midrule
$\theta$ & \descTheta & $tCO_2e$ \\
\bottomrule
\end{tabular}}

% ---------------------------------------------------------------------
% 5. Parámetros que el modelo conductual agrega a Amigo et al. (2021)
% ---------------------------------------------------------------------
\newcommand{\tablaParametrosAgregados}{%
\begin{tabular}{lll}
\toprule
\textbf{Símbolo} & \textbf{Descripción} & \textbf{Unidad} \\
\midrule
$\kappa$                & \descKappa      & \unidadUSDt \\
$CAP_s$                 & \descCAPs       & tCO$_2$e \\
$CAP_d$                 & \descCAPd       & tCO$_2$e \\
$\sigma_\varepsilon^2$  & \descSigmaEps   & tCO$_2$e \\
$\sigma_{CAP}^2$        & \descSigmaCAP   & tCO$_2$e \\
\bottomrule
\end{tabular}}

% ---------------------------------------------------------------------
% 6. Valores base de los parámetros, con su origen
% ---------------------------------------------------------------------
% Las descripciones de origen se definen aparte para que la
% presentación pueda usarlas en tarjetas.
\newcommand{\origenCAPs}{\refcita{ministerio_de_energia_de_chile_plan_2024}}
\newcommand{\origenCAPd}{Promedio sector energía 2012--2022 (SNICHILE)}
\newcommand{\origenSigmaCAP}{Incertidumbre estructural histórica}
\newcommand{\origenKappa}{Costo social del carbono, Gobierno de Chile (2024)}

\newcommand{\tablaParametrosChile}{%
\begin{tabular}{l c l l}
\toprule
\textbf{Parámetro} & \textbf{Valor} & \textbf{Unidad} & \textbf{Origen} \\
\midrule
$CAP_s$          & \valCAPs      & \unidadMt   & \origenCAPs \\
$CAP_d$          & \valCAPd      & \unidadMt   & \origenCAPd \\
$\sigma_{CAP}$   & \valSigmaCAP  & \unidadMt   & \origenSigmaCAP \\
$\kappa$         & \valKappa     & \unidadUSDt & \origenKappa \\
\bottomrule
\end{tabular}}

%
% Los macros de tabla entregan sólo el `tabular`, sin flotante ni
% caption, para que cada documento lo envuelva a su manera (`table` en
% la memoria, `frame` en beamer).
%
% Requiere: amsmath, booktabs.
% =====================================================================

\providecommand{\refcita}[1]{\texttt{#1}}

% ---------------------------------------------------------------------
% 1. Valores de los parámetros calibrados para Chile
% ---------------------------------------------------------------------
% Valores tal como los reporta el cuadro de parámetros base de la
% memoria. \valCAPdExacto es el valor sin redondear que reproducen las
% tablas de resultados; el capítulo de conclusiones usa 311,01.
% Pendiente de unificar por los autores: ver
% docs/auditoria/AUDITORIA_2_MHurtado.md.
\newcommand{\valCAPs}{264,20}          % MtCO2e
\newcommand{\valCAPd}{311,00}          % MtCO2e
\newcommand{\valCAPdExacto}{311,0134}  % MtCO2e
\newcommand{\valSigmaCAP}{61,86}       % MtCO2e
\newcommand{\valKappa}{64,40}          % USD/tCO2e
% Calibración: hay que distinguir dos valores de m que la memoria usa para
% cosas distintas (chapter04.tex, Cuadro tab:top10_distancia_mix y el texto
% que le sigue).
%   - El argmin de la distancia absoluta es m=0, que exige sigma_eps -> infinito
%     y por eso no sirve como parámetro operativo.
%   - Para los análisis posteriores se selecciona m=0,1, el menor m con un
%     sigma_eps finito, cuya distancia (0,2758) queda muy cerca del mínimo.
% Nota: el texto de la memoria dice sigma_eps = 186,90 y el cuadro dice 185,58
% para el mismo m=0,1. Pendiente de unificar por los autores; ver
% docs/auditoria/AUDITORIA_2_MHurtado.md.
\newcommand{\valMArgmin}{0}                  % menor distancia absoluta
\newcommand{\valDistanciaArgmin}{0,2727}
\newcommand{\valMSeleccionado}{0,1}          % usado en los análisis posteriores
\newcommand{\valDistanciaSeleccionado}{0,2758}
\newcommand{\valSigmaEpsSeleccionado}{185,58} % MtCO2e, segun tab:top10_distancia_mix
\newcommand{\valHorizonte}{2020--2030}
\newcommand{\unidadMt}{MtCO$_2$e}
\newcommand{\unidadUSDt}{USD/tCO$_2$e}

% ---------------------------------------------------------------------
% 2. Ecuaciones del planificador con inatención conductual
% ---------------------------------------------------------------------
% Cuerpo matemático sin entorno, para poder usarlo en `equation`,
% en \[...\] o dentro de un nodo tikz.
\newcommand{\ecSenal}{CAP_s = CAP + \epsilon}
\newcommand{\ecEsperanza}{\mathbb{E}\left[CAP|CAP_s\right] = m CAP_s + (1 - m)CAP_d}
\newcommand{\ecAtencion}{m = \frac{\sigma_{CAP}^2}{\sigma_{CAP}^2 + \sigma_\varepsilon^2}}
\newcommand{\ecThetaSinPrecio}{{\theta^*}(CAP_s) = m\cdot CAP_s + (1 - m)\cdot CAP_d}
\newcommand{\ecUtilidadPlanificador}{\mathcal{V}(\theta,CAP) = \pi^a \theta - \frac{\kappa}{2}(\theta - CAP)^2}
\newcommand{\ecThetaOptimo}{\theta^*(CAP_s) = m\,CAP_s + (1-m)\,CAP_d + \frac{\pi^a}{\kappa}}

% ---------------------------------------------------------------------
% 3. Descripciones de los símbolos, definidas una sola vez
% ---------------------------------------------------------------------
% La memoria las usa como filas de tabla y la presentación como
% tarjetas; por eso el texto se define aquí y el formato en cada
% documento.
\newcommand{\descCAP}{Presupuesto ambiental real}
\newcommand{\descCAPs}{Señal observada del presupuesto de emisiones}
\newcommand{\descCAPd}{Presupuesto determinista del CAP (valor por defecto)}
\newcommand{\descEps}{Ruido en la señal}
\newcommand{\descSigmaEps}{Varianza del error en la señal}
\newcommand{\descSigmaCAP}{Varianza del CAP}
\newcommand{\descTheta}{Permisos totales emitidos por el planificador}
\newcommand{\descPiA}{Precio de permisos emitidos por el planificador}
\newcommand{\descPiV}{Precio de permisos en el mercado entre productores}
\newcommand{\descPiD}{Precio de la electricidad pagado por el consumidor}
\newcommand{\descKappa}{Costo social por alejarse del $CAP$}
\newcommand{\descM}{Nivel de atención del planificador}
% Distribuciones
\newcommand{\distCAP}{CAP \sim \mathcal{N}(CAP_d,\sigma^2_{CAP})}
\newcommand{\distEps}{\epsilon \sim \mathcal{N}(0,\sigma^2_\varepsilon)}

% ---------------------------------------------------------------------
% 4. Nomenclatura del modelo de Amigo et al. (2021)
% ---------------------------------------------------------------------
\newcommand{\tablaNomenclaturaAmigo}{%
\begin{tabular}{lll}
\toprule
\textbf{Símbolo} & \textbf{Descripción} & \textbf{Unidad} \\
\midrule
$\pi^{a}$ & \descPiA & USD/$tCO_2e$ \\
$\pi^{v}$ & \descPiV & USD/permiso \\
$\pi^{d}$ & \descPiD & USD/MWh \\
\midrule
$x_i$ & Inversión o expansión de capacidad del productor $i$ & MW \\
$Q_i$ & Generación eléctrica del productor $i$ & MWh \\
$A_i$ & Permisos adquiridos por el productor $i$ al planificador & $tCO_2e$\\
$P_i$ & Permisos comprados por el productor $i$ en el mercado secundario & $tCO_2e$ \\
$V_i$ & Permisos vendidos por el productor $i$ en el mercado secundario & $tCO_2e$ \\
$\epsilon_i$ & Factor de emisión por tecnología $i$ & $tCO_2$e/MWh \\
\midrule
$\theta$ & \descTheta & $tCO_2e$ \\
\bottomrule
\end{tabular}}

% ---------------------------------------------------------------------
% 5. Parámetros que el modelo conductual agrega a Amigo et al. (2021)
% ---------------------------------------------------------------------
\newcommand{\tablaParametrosAgregados}{%
\begin{tabular}{lll}
\toprule
\textbf{Símbolo} & \textbf{Descripción} & \textbf{Unidad} \\
\midrule
$\kappa$                & \descKappa      & \unidadUSDt \\
$CAP_s$                 & \descCAPs       & tCO$_2$e \\
$CAP_d$                 & \descCAPd       & tCO$_2$e \\
$\sigma_\varepsilon^2$  & \descSigmaEps   & tCO$_2$e \\
$\sigma_{CAP}^2$        & \descSigmaCAP   & tCO$_2$e \\
\bottomrule
\end{tabular}}

% ---------------------------------------------------------------------
% 6. Valores base de los parámetros, con su origen
% ---------------------------------------------------------------------
% Las descripciones de origen se definen aparte para que la
% presentación pueda usarlas en tarjetas.
\newcommand{\origenCAPs}{\refcita{ministerio_de_energia_de_chile_plan_2024}}
\newcommand{\origenCAPd}{Promedio sector energía 2012--2022 (SNICHILE)}
\newcommand{\origenSigmaCAP}{Incertidumbre estructural histórica}
\newcommand{\origenKappa}{Costo social del carbono, Gobierno de Chile (2024)}

\newcommand{\tablaParametrosChile}{%
\begin{tabular}{l c l l}
\toprule
\textbf{Parámetro} & \textbf{Valor} & \textbf{Unidad} & \textbf{Origen} \\
\midrule
$CAP_s$          & \valCAPs      & \unidadMt   & \origenCAPs \\
$CAP_d$          & \valCAPd      & \unidadMt   & \origenCAPd \\
$\sigma_{CAP}$   & \valSigmaCAP  & \unidadMt   & \origenSigmaCAP \\
$\kappa$         & \valKappa     & \unidadUSDt & \origenKappa \\
\bottomrule
\end{tabular}}

% o
%   \newcommand{\refcita}[1]{\citet{#1}}   % presentación (natbib)
%   % =====================================================================
% shared/datos-modelo.tex
% ---------------------------------------------------------------------
% Nomenclatura, ecuaciones y valores de parámetros del modelo de
% anclaje y ajuste del subastador. Fuente única para la memoria
% (docs/MemoriaMHurtado) y la presentación de defensa
% (docs/Presentations/presentation-template), de modo que un valor o una
% ecuación se escriban una sola vez.
%
% Uso:
%   \newcommand{\refcita}[1]{\citeA{#1}}   % memoria (newapa/apacite)
%   % =====================================================================
% shared/datos-modelo.tex
% ---------------------------------------------------------------------
% Nomenclatura, ecuaciones y valores de parámetros del modelo de
% anclaje y ajuste del subastador. Fuente única para la memoria
% (docs/MemoriaMHurtado) y la presentación de defensa
% (docs/Presentations/presentation-template), de modo que un valor o una
% ecuación se escriban una sola vez.
%
% Uso:
%   \newcommand{\refcita}[1]{\citeA{#1}}   % memoria (newapa/apacite)
%   % =====================================================================
% shared/datos-modelo.tex
% ---------------------------------------------------------------------
% Nomenclatura, ecuaciones y valores de parámetros del modelo de
% anclaje y ajuste del subastador. Fuente única para la memoria
% (docs/MemoriaMHurtado) y la presentación de defensa
% (docs/Presentations/presentation-template), de modo que un valor o una
% ecuación se escriban una sola vez.
%
% Uso:
%   \newcommand{\refcita}[1]{\citeA{#1}}   % memoria (newapa/apacite)
%   \input{../../shared/datos-modelo}
% o
%   \newcommand{\refcita}[1]{\citet{#1}}   % presentación (natbib)
%   \input{../../../shared/datos-modelo}
%
% Los macros de tabla entregan sólo el `tabular`, sin flotante ni
% caption, para que cada documento lo envuelva a su manera (`table` en
% la memoria, `frame` en beamer).
%
% Requiere: amsmath, booktabs.
% =====================================================================

\providecommand{\refcita}[1]{\texttt{#1}}

% ---------------------------------------------------------------------
% 1. Valores de los parámetros calibrados para Chile
% ---------------------------------------------------------------------
% Valores tal como los reporta el cuadro de parámetros base de la
% memoria. \valCAPdExacto es el valor sin redondear que reproducen las
% tablas de resultados; el capítulo de conclusiones usa 311,01.
% Pendiente de unificar por los autores: ver
% docs/auditoria/AUDITORIA_2_MHurtado.md.
\newcommand{\valCAPs}{264,20}          % MtCO2e
\newcommand{\valCAPd}{311,00}          % MtCO2e
\newcommand{\valCAPdExacto}{311,0134}  % MtCO2e
\newcommand{\valSigmaCAP}{61,86}       % MtCO2e
\newcommand{\valKappa}{64,40}          % USD/tCO2e
% Calibración: hay que distinguir dos valores de m que la memoria usa para
% cosas distintas (chapter04.tex, Cuadro tab:top10_distancia_mix y el texto
% que le sigue).
%   - El argmin de la distancia absoluta es m=0, que exige sigma_eps -> infinito
%     y por eso no sirve como parámetro operativo.
%   - Para los análisis posteriores se selecciona m=0,1, el menor m con un
%     sigma_eps finito, cuya distancia (0,2758) queda muy cerca del mínimo.
% Nota: el texto de la memoria dice sigma_eps = 186,90 y el cuadro dice 185,58
% para el mismo m=0,1. Pendiente de unificar por los autores; ver
% docs/auditoria/AUDITORIA_2_MHurtado.md.
\newcommand{\valMArgmin}{0}                  % menor distancia absoluta
\newcommand{\valDistanciaArgmin}{0,2727}
\newcommand{\valMSeleccionado}{0,1}          % usado en los análisis posteriores
\newcommand{\valDistanciaSeleccionado}{0,2758}
\newcommand{\valSigmaEpsSeleccionado}{185,58} % MtCO2e, segun tab:top10_distancia_mix
\newcommand{\valHorizonte}{2020--2030}
\newcommand{\unidadMt}{MtCO$_2$e}
\newcommand{\unidadUSDt}{USD/tCO$_2$e}

% ---------------------------------------------------------------------
% 2. Ecuaciones del planificador con inatención conductual
% ---------------------------------------------------------------------
% Cuerpo matemático sin entorno, para poder usarlo en `equation`,
% en \[...\] o dentro de un nodo tikz.
\newcommand{\ecSenal}{CAP_s = CAP + \epsilon}
\newcommand{\ecEsperanza}{\mathbb{E}\left[CAP|CAP_s\right] = m CAP_s + (1 - m)CAP_d}
\newcommand{\ecAtencion}{m = \frac{\sigma_{CAP}^2}{\sigma_{CAP}^2 + \sigma_\varepsilon^2}}
\newcommand{\ecThetaSinPrecio}{{\theta^*}(CAP_s) = m\cdot CAP_s + (1 - m)\cdot CAP_d}
\newcommand{\ecUtilidadPlanificador}{\mathcal{V}(\theta,CAP) = \pi^a \theta - \frac{\kappa}{2}(\theta - CAP)^2}
\newcommand{\ecThetaOptimo}{\theta^*(CAP_s) = m\,CAP_s + (1-m)\,CAP_d + \frac{\pi^a}{\kappa}}

% ---------------------------------------------------------------------
% 3. Descripciones de los símbolos, definidas una sola vez
% ---------------------------------------------------------------------
% La memoria las usa como filas de tabla y la presentación como
% tarjetas; por eso el texto se define aquí y el formato en cada
% documento.
\newcommand{\descCAP}{Presupuesto ambiental real}
\newcommand{\descCAPs}{Señal observada del presupuesto de emisiones}
\newcommand{\descCAPd}{Presupuesto determinista del CAP (valor por defecto)}
\newcommand{\descEps}{Ruido en la señal}
\newcommand{\descSigmaEps}{Varianza del error en la señal}
\newcommand{\descSigmaCAP}{Varianza del CAP}
\newcommand{\descTheta}{Permisos totales emitidos por el planificador}
\newcommand{\descPiA}{Precio de permisos emitidos por el planificador}
\newcommand{\descPiV}{Precio de permisos en el mercado entre productores}
\newcommand{\descPiD}{Precio de la electricidad pagado por el consumidor}
\newcommand{\descKappa}{Costo social por alejarse del $CAP$}
\newcommand{\descM}{Nivel de atención del planificador}
% Distribuciones
\newcommand{\distCAP}{CAP \sim \mathcal{N}(CAP_d,\sigma^2_{CAP})}
\newcommand{\distEps}{\epsilon \sim \mathcal{N}(0,\sigma^2_\varepsilon)}

% ---------------------------------------------------------------------
% 4. Nomenclatura del modelo de Amigo et al. (2021)
% ---------------------------------------------------------------------
\newcommand{\tablaNomenclaturaAmigo}{%
\begin{tabular}{lll}
\toprule
\textbf{Símbolo} & \textbf{Descripción} & \textbf{Unidad} \\
\midrule
$\pi^{a}$ & \descPiA & USD/$tCO_2e$ \\
$\pi^{v}$ & \descPiV & USD/permiso \\
$\pi^{d}$ & \descPiD & USD/MWh \\
\midrule
$x_i$ & Inversión o expansión de capacidad del productor $i$ & MW \\
$Q_i$ & Generación eléctrica del productor $i$ & MWh \\
$A_i$ & Permisos adquiridos por el productor $i$ al planificador & $tCO_2e$\\
$P_i$ & Permisos comprados por el productor $i$ en el mercado secundario & $tCO_2e$ \\
$V_i$ & Permisos vendidos por el productor $i$ en el mercado secundario & $tCO_2e$ \\
$\epsilon_i$ & Factor de emisión por tecnología $i$ & $tCO_2$e/MWh \\
\midrule
$\theta$ & \descTheta & $tCO_2e$ \\
\bottomrule
\end{tabular}}

% ---------------------------------------------------------------------
% 5. Parámetros que el modelo conductual agrega a Amigo et al. (2021)
% ---------------------------------------------------------------------
\newcommand{\tablaParametrosAgregados}{%
\begin{tabular}{lll}
\toprule
\textbf{Símbolo} & \textbf{Descripción} & \textbf{Unidad} \\
\midrule
$\kappa$                & \descKappa      & \unidadUSDt \\
$CAP_s$                 & \descCAPs       & tCO$_2$e \\
$CAP_d$                 & \descCAPd       & tCO$_2$e \\
$\sigma_\varepsilon^2$  & \descSigmaEps   & tCO$_2$e \\
$\sigma_{CAP}^2$        & \descSigmaCAP   & tCO$_2$e \\
\bottomrule
\end{tabular}}

% ---------------------------------------------------------------------
% 6. Valores base de los parámetros, con su origen
% ---------------------------------------------------------------------
% Las descripciones de origen se definen aparte para que la
% presentación pueda usarlas en tarjetas.
\newcommand{\origenCAPs}{\refcita{ministerio_de_energia_de_chile_plan_2024}}
\newcommand{\origenCAPd}{Promedio sector energía 2012--2022 (SNICHILE)}
\newcommand{\origenSigmaCAP}{Incertidumbre estructural histórica}
\newcommand{\origenKappa}{Costo social del carbono, Gobierno de Chile (2024)}

\newcommand{\tablaParametrosChile}{%
\begin{tabular}{l c l l}
\toprule
\textbf{Parámetro} & \textbf{Valor} & \textbf{Unidad} & \textbf{Origen} \\
\midrule
$CAP_s$          & \valCAPs      & \unidadMt   & \origenCAPs \\
$CAP_d$          & \valCAPd      & \unidadMt   & \origenCAPd \\
$\sigma_{CAP}$   & \valSigmaCAP  & \unidadMt   & \origenSigmaCAP \\
$\kappa$         & \valKappa     & \unidadUSDt & \origenKappa \\
\bottomrule
\end{tabular}}

% o
%   \newcommand{\refcita}[1]{\citet{#1}}   % presentación (natbib)
%   % =====================================================================
% shared/datos-modelo.tex
% ---------------------------------------------------------------------
% Nomenclatura, ecuaciones y valores de parámetros del modelo de
% anclaje y ajuste del subastador. Fuente única para la memoria
% (docs/MemoriaMHurtado) y la presentación de defensa
% (docs/Presentations/presentation-template), de modo que un valor o una
% ecuación se escriban una sola vez.
%
% Uso:
%   \newcommand{\refcita}[1]{\citeA{#1}}   % memoria (newapa/apacite)
%   \input{../../shared/datos-modelo}
% o
%   \newcommand{\refcita}[1]{\citet{#1}}   % presentación (natbib)
%   \input{../../../shared/datos-modelo}
%
% Los macros de tabla entregan sólo el `tabular`, sin flotante ni
% caption, para que cada documento lo envuelva a su manera (`table` en
% la memoria, `frame` en beamer).
%
% Requiere: amsmath, booktabs.
% =====================================================================

\providecommand{\refcita}[1]{\texttt{#1}}

% ---------------------------------------------------------------------
% 1. Valores de los parámetros calibrados para Chile
% ---------------------------------------------------------------------
% Valores tal como los reporta el cuadro de parámetros base de la
% memoria. \valCAPdExacto es el valor sin redondear que reproducen las
% tablas de resultados; el capítulo de conclusiones usa 311,01.
% Pendiente de unificar por los autores: ver
% docs/auditoria/AUDITORIA_2_MHurtado.md.
\newcommand{\valCAPs}{264,20}          % MtCO2e
\newcommand{\valCAPd}{311,00}          % MtCO2e
\newcommand{\valCAPdExacto}{311,0134}  % MtCO2e
\newcommand{\valSigmaCAP}{61,86}       % MtCO2e
\newcommand{\valKappa}{64,40}          % USD/tCO2e
% Calibración: hay que distinguir dos valores de m que la memoria usa para
% cosas distintas (chapter04.tex, Cuadro tab:top10_distancia_mix y el texto
% que le sigue).
%   - El argmin de la distancia absoluta es m=0, que exige sigma_eps -> infinito
%     y por eso no sirve como parámetro operativo.
%   - Para los análisis posteriores se selecciona m=0,1, el menor m con un
%     sigma_eps finito, cuya distancia (0,2758) queda muy cerca del mínimo.
% Nota: el texto de la memoria dice sigma_eps = 186,90 y el cuadro dice 185,58
% para el mismo m=0,1. Pendiente de unificar por los autores; ver
% docs/auditoria/AUDITORIA_2_MHurtado.md.
\newcommand{\valMArgmin}{0}                  % menor distancia absoluta
\newcommand{\valDistanciaArgmin}{0,2727}
\newcommand{\valMSeleccionado}{0,1}          % usado en los análisis posteriores
\newcommand{\valDistanciaSeleccionado}{0,2758}
\newcommand{\valSigmaEpsSeleccionado}{185,58} % MtCO2e, segun tab:top10_distancia_mix
\newcommand{\valHorizonte}{2020--2030}
\newcommand{\unidadMt}{MtCO$_2$e}
\newcommand{\unidadUSDt}{USD/tCO$_2$e}

% ---------------------------------------------------------------------
% 2. Ecuaciones del planificador con inatención conductual
% ---------------------------------------------------------------------
% Cuerpo matemático sin entorno, para poder usarlo en `equation`,
% en \[...\] o dentro de un nodo tikz.
\newcommand{\ecSenal}{CAP_s = CAP + \epsilon}
\newcommand{\ecEsperanza}{\mathbb{E}\left[CAP|CAP_s\right] = m CAP_s + (1 - m)CAP_d}
\newcommand{\ecAtencion}{m = \frac{\sigma_{CAP}^2}{\sigma_{CAP}^2 + \sigma_\varepsilon^2}}
\newcommand{\ecThetaSinPrecio}{{\theta^*}(CAP_s) = m\cdot CAP_s + (1 - m)\cdot CAP_d}
\newcommand{\ecUtilidadPlanificador}{\mathcal{V}(\theta,CAP) = \pi^a \theta - \frac{\kappa}{2}(\theta - CAP)^2}
\newcommand{\ecThetaOptimo}{\theta^*(CAP_s) = m\,CAP_s + (1-m)\,CAP_d + \frac{\pi^a}{\kappa}}

% ---------------------------------------------------------------------
% 3. Descripciones de los símbolos, definidas una sola vez
% ---------------------------------------------------------------------
% La memoria las usa como filas de tabla y la presentación como
% tarjetas; por eso el texto se define aquí y el formato en cada
% documento.
\newcommand{\descCAP}{Presupuesto ambiental real}
\newcommand{\descCAPs}{Señal observada del presupuesto de emisiones}
\newcommand{\descCAPd}{Presupuesto determinista del CAP (valor por defecto)}
\newcommand{\descEps}{Ruido en la señal}
\newcommand{\descSigmaEps}{Varianza del error en la señal}
\newcommand{\descSigmaCAP}{Varianza del CAP}
\newcommand{\descTheta}{Permisos totales emitidos por el planificador}
\newcommand{\descPiA}{Precio de permisos emitidos por el planificador}
\newcommand{\descPiV}{Precio de permisos en el mercado entre productores}
\newcommand{\descPiD}{Precio de la electricidad pagado por el consumidor}
\newcommand{\descKappa}{Costo social por alejarse del $CAP$}
\newcommand{\descM}{Nivel de atención del planificador}
% Distribuciones
\newcommand{\distCAP}{CAP \sim \mathcal{N}(CAP_d,\sigma^2_{CAP})}
\newcommand{\distEps}{\epsilon \sim \mathcal{N}(0,\sigma^2_\varepsilon)}

% ---------------------------------------------------------------------
% 4. Nomenclatura del modelo de Amigo et al. (2021)
% ---------------------------------------------------------------------
\newcommand{\tablaNomenclaturaAmigo}{%
\begin{tabular}{lll}
\toprule
\textbf{Símbolo} & \textbf{Descripción} & \textbf{Unidad} \\
\midrule
$\pi^{a}$ & \descPiA & USD/$tCO_2e$ \\
$\pi^{v}$ & \descPiV & USD/permiso \\
$\pi^{d}$ & \descPiD & USD/MWh \\
\midrule
$x_i$ & Inversión o expansión de capacidad del productor $i$ & MW \\
$Q_i$ & Generación eléctrica del productor $i$ & MWh \\
$A_i$ & Permisos adquiridos por el productor $i$ al planificador & $tCO_2e$\\
$P_i$ & Permisos comprados por el productor $i$ en el mercado secundario & $tCO_2e$ \\
$V_i$ & Permisos vendidos por el productor $i$ en el mercado secundario & $tCO_2e$ \\
$\epsilon_i$ & Factor de emisión por tecnología $i$ & $tCO_2$e/MWh \\
\midrule
$\theta$ & \descTheta & $tCO_2e$ \\
\bottomrule
\end{tabular}}

% ---------------------------------------------------------------------
% 5. Parámetros que el modelo conductual agrega a Amigo et al. (2021)
% ---------------------------------------------------------------------
\newcommand{\tablaParametrosAgregados}{%
\begin{tabular}{lll}
\toprule
\textbf{Símbolo} & \textbf{Descripción} & \textbf{Unidad} \\
\midrule
$\kappa$                & \descKappa      & \unidadUSDt \\
$CAP_s$                 & \descCAPs       & tCO$_2$e \\
$CAP_d$                 & \descCAPd       & tCO$_2$e \\
$\sigma_\varepsilon^2$  & \descSigmaEps   & tCO$_2$e \\
$\sigma_{CAP}^2$        & \descSigmaCAP   & tCO$_2$e \\
\bottomrule
\end{tabular}}

% ---------------------------------------------------------------------
% 6. Valores base de los parámetros, con su origen
% ---------------------------------------------------------------------
% Las descripciones de origen se definen aparte para que la
% presentación pueda usarlas en tarjetas.
\newcommand{\origenCAPs}{\refcita{ministerio_de_energia_de_chile_plan_2024}}
\newcommand{\origenCAPd}{Promedio sector energía 2012--2022 (SNICHILE)}
\newcommand{\origenSigmaCAP}{Incertidumbre estructural histórica}
\newcommand{\origenKappa}{Costo social del carbono, Gobierno de Chile (2024)}

\newcommand{\tablaParametrosChile}{%
\begin{tabular}{l c l l}
\toprule
\textbf{Parámetro} & \textbf{Valor} & \textbf{Unidad} & \textbf{Origen} \\
\midrule
$CAP_s$          & \valCAPs      & \unidadMt   & \origenCAPs \\
$CAP_d$          & \valCAPd      & \unidadMt   & \origenCAPd \\
$\sigma_{CAP}$   & \valSigmaCAP  & \unidadMt   & \origenSigmaCAP \\
$\kappa$         & \valKappa     & \unidadUSDt & \origenKappa \\
\bottomrule
\end{tabular}}

%
% Los macros de tabla entregan sólo el `tabular`, sin flotante ni
% caption, para que cada documento lo envuelva a su manera (`table` en
% la memoria, `frame` en beamer).
%
% Requiere: amsmath, booktabs.
% =====================================================================

\providecommand{\refcita}[1]{\texttt{#1}}

% ---------------------------------------------------------------------
% 1. Valores de los parámetros calibrados para Chile
% ---------------------------------------------------------------------
% Valores tal como los reporta el cuadro de parámetros base de la
% memoria. \valCAPdExacto es el valor sin redondear que reproducen las
% tablas de resultados; el capítulo de conclusiones usa 311,01.
% Pendiente de unificar por los autores: ver
% docs/auditoria/AUDITORIA_2_MHurtado.md.
\newcommand{\valCAPs}{264,20}          % MtCO2e
\newcommand{\valCAPd}{311,00}          % MtCO2e
\newcommand{\valCAPdExacto}{311,0134}  % MtCO2e
\newcommand{\valSigmaCAP}{61,86}       % MtCO2e
\newcommand{\valKappa}{64,40}          % USD/tCO2e
% Calibración: hay que distinguir dos valores de m que la memoria usa para
% cosas distintas (chapter04.tex, Cuadro tab:top10_distancia_mix y el texto
% que le sigue).
%   - El argmin de la distancia absoluta es m=0, que exige sigma_eps -> infinito
%     y por eso no sirve como parámetro operativo.
%   - Para los análisis posteriores se selecciona m=0,1, el menor m con un
%     sigma_eps finito, cuya distancia (0,2758) queda muy cerca del mínimo.
% Nota: el texto de la memoria dice sigma_eps = 186,90 y el cuadro dice 185,58
% para el mismo m=0,1. Pendiente de unificar por los autores; ver
% docs/auditoria/AUDITORIA_2_MHurtado.md.
\newcommand{\valMArgmin}{0}                  % menor distancia absoluta
\newcommand{\valDistanciaArgmin}{0,2727}
\newcommand{\valMSeleccionado}{0,1}          % usado en los análisis posteriores
\newcommand{\valDistanciaSeleccionado}{0,2758}
\newcommand{\valSigmaEpsSeleccionado}{185,58} % MtCO2e, segun tab:top10_distancia_mix
\newcommand{\valHorizonte}{2020--2030}
\newcommand{\unidadMt}{MtCO$_2$e}
\newcommand{\unidadUSDt}{USD/tCO$_2$e}

% ---------------------------------------------------------------------
% 2. Ecuaciones del planificador con inatención conductual
% ---------------------------------------------------------------------
% Cuerpo matemático sin entorno, para poder usarlo en `equation`,
% en \[...\] o dentro de un nodo tikz.
\newcommand{\ecSenal}{CAP_s = CAP + \epsilon}
\newcommand{\ecEsperanza}{\mathbb{E}\left[CAP|CAP_s\right] = m CAP_s + (1 - m)CAP_d}
\newcommand{\ecAtencion}{m = \frac{\sigma_{CAP}^2}{\sigma_{CAP}^2 + \sigma_\varepsilon^2}}
\newcommand{\ecThetaSinPrecio}{{\theta^*}(CAP_s) = m\cdot CAP_s + (1 - m)\cdot CAP_d}
\newcommand{\ecUtilidadPlanificador}{\mathcal{V}(\theta,CAP) = \pi^a \theta - \frac{\kappa}{2}(\theta - CAP)^2}
\newcommand{\ecThetaOptimo}{\theta^*(CAP_s) = m\,CAP_s + (1-m)\,CAP_d + \frac{\pi^a}{\kappa}}

% ---------------------------------------------------------------------
% 3. Descripciones de los símbolos, definidas una sola vez
% ---------------------------------------------------------------------
% La memoria las usa como filas de tabla y la presentación como
% tarjetas; por eso el texto se define aquí y el formato en cada
% documento.
\newcommand{\descCAP}{Presupuesto ambiental real}
\newcommand{\descCAPs}{Señal observada del presupuesto de emisiones}
\newcommand{\descCAPd}{Presupuesto determinista del CAP (valor por defecto)}
\newcommand{\descEps}{Ruido en la señal}
\newcommand{\descSigmaEps}{Varianza del error en la señal}
\newcommand{\descSigmaCAP}{Varianza del CAP}
\newcommand{\descTheta}{Permisos totales emitidos por el planificador}
\newcommand{\descPiA}{Precio de permisos emitidos por el planificador}
\newcommand{\descPiV}{Precio de permisos en el mercado entre productores}
\newcommand{\descPiD}{Precio de la electricidad pagado por el consumidor}
\newcommand{\descKappa}{Costo social por alejarse del $CAP$}
\newcommand{\descM}{Nivel de atención del planificador}
% Distribuciones
\newcommand{\distCAP}{CAP \sim \mathcal{N}(CAP_d,\sigma^2_{CAP})}
\newcommand{\distEps}{\epsilon \sim \mathcal{N}(0,\sigma^2_\varepsilon)}

% ---------------------------------------------------------------------
% 4. Nomenclatura del modelo de Amigo et al. (2021)
% ---------------------------------------------------------------------
\newcommand{\tablaNomenclaturaAmigo}{%
\begin{tabular}{lll}
\toprule
\textbf{Símbolo} & \textbf{Descripción} & \textbf{Unidad} \\
\midrule
$\pi^{a}$ & \descPiA & USD/$tCO_2e$ \\
$\pi^{v}$ & \descPiV & USD/permiso \\
$\pi^{d}$ & \descPiD & USD/MWh \\
\midrule
$x_i$ & Inversión o expansión de capacidad del productor $i$ & MW \\
$Q_i$ & Generación eléctrica del productor $i$ & MWh \\
$A_i$ & Permisos adquiridos por el productor $i$ al planificador & $tCO_2e$\\
$P_i$ & Permisos comprados por el productor $i$ en el mercado secundario & $tCO_2e$ \\
$V_i$ & Permisos vendidos por el productor $i$ en el mercado secundario & $tCO_2e$ \\
$\epsilon_i$ & Factor de emisión por tecnología $i$ & $tCO_2$e/MWh \\
\midrule
$\theta$ & \descTheta & $tCO_2e$ \\
\bottomrule
\end{tabular}}

% ---------------------------------------------------------------------
% 5. Parámetros que el modelo conductual agrega a Amigo et al. (2021)
% ---------------------------------------------------------------------
\newcommand{\tablaParametrosAgregados}{%
\begin{tabular}{lll}
\toprule
\textbf{Símbolo} & \textbf{Descripción} & \textbf{Unidad} \\
\midrule
$\kappa$                & \descKappa      & \unidadUSDt \\
$CAP_s$                 & \descCAPs       & tCO$_2$e \\
$CAP_d$                 & \descCAPd       & tCO$_2$e \\
$\sigma_\varepsilon^2$  & \descSigmaEps   & tCO$_2$e \\
$\sigma_{CAP}^2$        & \descSigmaCAP   & tCO$_2$e \\
\bottomrule
\end{tabular}}

% ---------------------------------------------------------------------
% 6. Valores base de los parámetros, con su origen
% ---------------------------------------------------------------------
% Las descripciones de origen se definen aparte para que la
% presentación pueda usarlas en tarjetas.
\newcommand{\origenCAPs}{\refcita{ministerio_de_energia_de_chile_plan_2024}}
\newcommand{\origenCAPd}{Promedio sector energía 2012--2022 (SNICHILE)}
\newcommand{\origenSigmaCAP}{Incertidumbre estructural histórica}
\newcommand{\origenKappa}{Costo social del carbono, Gobierno de Chile (2024)}

\newcommand{\tablaParametrosChile}{%
\begin{tabular}{l c l l}
\toprule
\textbf{Parámetro} & \textbf{Valor} & \textbf{Unidad} & \textbf{Origen} \\
\midrule
$CAP_s$          & \valCAPs      & \unidadMt   & \origenCAPs \\
$CAP_d$          & \valCAPd      & \unidadMt   & \origenCAPd \\
$\sigma_{CAP}$   & \valSigmaCAP  & \unidadMt   & \origenSigmaCAP \\
$\kappa$         & \valKappa     & \unidadUSDt & \origenKappa \\
\bottomrule
\end{tabular}}

% o
%   \newcommand{\refcita}[1]{\citet{#1}}   % presentación (natbib)
%   % =====================================================================
% shared/datos-modelo.tex
% ---------------------------------------------------------------------
% Nomenclatura, ecuaciones y valores de parámetros del modelo de
% anclaje y ajuste del subastador. Fuente única para la memoria
% (docs/MemoriaMHurtado) y la presentación de defensa
% (docs/Presentations/presentation-template), de modo que un valor o una
% ecuación se escriban una sola vez.
%
% Uso:
%   \newcommand{\refcita}[1]{\citeA{#1}}   % memoria (newapa/apacite)
%   % =====================================================================
% shared/datos-modelo.tex
% ---------------------------------------------------------------------
% Nomenclatura, ecuaciones y valores de parámetros del modelo de
% anclaje y ajuste del subastador. Fuente única para la memoria
% (docs/MemoriaMHurtado) y la presentación de defensa
% (docs/Presentations/presentation-template), de modo que un valor o una
% ecuación se escriban una sola vez.
%
% Uso:
%   \newcommand{\refcita}[1]{\citeA{#1}}   % memoria (newapa/apacite)
%   \input{../../shared/datos-modelo}
% o
%   \newcommand{\refcita}[1]{\citet{#1}}   % presentación (natbib)
%   \input{../../../shared/datos-modelo}
%
% Los macros de tabla entregan sólo el `tabular`, sin flotante ni
% caption, para que cada documento lo envuelva a su manera (`table` en
% la memoria, `frame` en beamer).
%
% Requiere: amsmath, booktabs.
% =====================================================================

\providecommand{\refcita}[1]{\texttt{#1}}

% ---------------------------------------------------------------------
% 1. Valores de los parámetros calibrados para Chile
% ---------------------------------------------------------------------
% Valores tal como los reporta el cuadro de parámetros base de la
% memoria. \valCAPdExacto es el valor sin redondear que reproducen las
% tablas de resultados; el capítulo de conclusiones usa 311,01.
% Pendiente de unificar por los autores: ver
% docs/auditoria/AUDITORIA_2_MHurtado.md.
\newcommand{\valCAPs}{264,20}          % MtCO2e
\newcommand{\valCAPd}{311,00}          % MtCO2e
\newcommand{\valCAPdExacto}{311,0134}  % MtCO2e
\newcommand{\valSigmaCAP}{61,86}       % MtCO2e
\newcommand{\valKappa}{64,40}          % USD/tCO2e
% Calibración: hay que distinguir dos valores de m que la memoria usa para
% cosas distintas (chapter04.tex, Cuadro tab:top10_distancia_mix y el texto
% que le sigue).
%   - El argmin de la distancia absoluta es m=0, que exige sigma_eps -> infinito
%     y por eso no sirve como parámetro operativo.
%   - Para los análisis posteriores se selecciona m=0,1, el menor m con un
%     sigma_eps finito, cuya distancia (0,2758) queda muy cerca del mínimo.
% Nota: el texto de la memoria dice sigma_eps = 186,90 y el cuadro dice 185,58
% para el mismo m=0,1. Pendiente de unificar por los autores; ver
% docs/auditoria/AUDITORIA_2_MHurtado.md.
\newcommand{\valMArgmin}{0}                  % menor distancia absoluta
\newcommand{\valDistanciaArgmin}{0,2727}
\newcommand{\valMSeleccionado}{0,1}          % usado en los análisis posteriores
\newcommand{\valDistanciaSeleccionado}{0,2758}
\newcommand{\valSigmaEpsSeleccionado}{185,58} % MtCO2e, segun tab:top10_distancia_mix
\newcommand{\valHorizonte}{2020--2030}
\newcommand{\unidadMt}{MtCO$_2$e}
\newcommand{\unidadUSDt}{USD/tCO$_2$e}

% ---------------------------------------------------------------------
% 2. Ecuaciones del planificador con inatención conductual
% ---------------------------------------------------------------------
% Cuerpo matemático sin entorno, para poder usarlo en `equation`,
% en \[...\] o dentro de un nodo tikz.
\newcommand{\ecSenal}{CAP_s = CAP + \epsilon}
\newcommand{\ecEsperanza}{\mathbb{E}\left[CAP|CAP_s\right] = m CAP_s + (1 - m)CAP_d}
\newcommand{\ecAtencion}{m = \frac{\sigma_{CAP}^2}{\sigma_{CAP}^2 + \sigma_\varepsilon^2}}
\newcommand{\ecThetaSinPrecio}{{\theta^*}(CAP_s) = m\cdot CAP_s + (1 - m)\cdot CAP_d}
\newcommand{\ecUtilidadPlanificador}{\mathcal{V}(\theta,CAP) = \pi^a \theta - \frac{\kappa}{2}(\theta - CAP)^2}
\newcommand{\ecThetaOptimo}{\theta^*(CAP_s) = m\,CAP_s + (1-m)\,CAP_d + \frac{\pi^a}{\kappa}}

% ---------------------------------------------------------------------
% 3. Descripciones de los símbolos, definidas una sola vez
% ---------------------------------------------------------------------
% La memoria las usa como filas de tabla y la presentación como
% tarjetas; por eso el texto se define aquí y el formato en cada
% documento.
\newcommand{\descCAP}{Presupuesto ambiental real}
\newcommand{\descCAPs}{Señal observada del presupuesto de emisiones}
\newcommand{\descCAPd}{Presupuesto determinista del CAP (valor por defecto)}
\newcommand{\descEps}{Ruido en la señal}
\newcommand{\descSigmaEps}{Varianza del error en la señal}
\newcommand{\descSigmaCAP}{Varianza del CAP}
\newcommand{\descTheta}{Permisos totales emitidos por el planificador}
\newcommand{\descPiA}{Precio de permisos emitidos por el planificador}
\newcommand{\descPiV}{Precio de permisos en el mercado entre productores}
\newcommand{\descPiD}{Precio de la electricidad pagado por el consumidor}
\newcommand{\descKappa}{Costo social por alejarse del $CAP$}
\newcommand{\descM}{Nivel de atención del planificador}
% Distribuciones
\newcommand{\distCAP}{CAP \sim \mathcal{N}(CAP_d,\sigma^2_{CAP})}
\newcommand{\distEps}{\epsilon \sim \mathcal{N}(0,\sigma^2_\varepsilon)}

% ---------------------------------------------------------------------
% 4. Nomenclatura del modelo de Amigo et al. (2021)
% ---------------------------------------------------------------------
\newcommand{\tablaNomenclaturaAmigo}{%
\begin{tabular}{lll}
\toprule
\textbf{Símbolo} & \textbf{Descripción} & \textbf{Unidad} \\
\midrule
$\pi^{a}$ & \descPiA & USD/$tCO_2e$ \\
$\pi^{v}$ & \descPiV & USD/permiso \\
$\pi^{d}$ & \descPiD & USD/MWh \\
\midrule
$x_i$ & Inversión o expansión de capacidad del productor $i$ & MW \\
$Q_i$ & Generación eléctrica del productor $i$ & MWh \\
$A_i$ & Permisos adquiridos por el productor $i$ al planificador & $tCO_2e$\\
$P_i$ & Permisos comprados por el productor $i$ en el mercado secundario & $tCO_2e$ \\
$V_i$ & Permisos vendidos por el productor $i$ en el mercado secundario & $tCO_2e$ \\
$\epsilon_i$ & Factor de emisión por tecnología $i$ & $tCO_2$e/MWh \\
\midrule
$\theta$ & \descTheta & $tCO_2e$ \\
\bottomrule
\end{tabular}}

% ---------------------------------------------------------------------
% 5. Parámetros que el modelo conductual agrega a Amigo et al. (2021)
% ---------------------------------------------------------------------
\newcommand{\tablaParametrosAgregados}{%
\begin{tabular}{lll}
\toprule
\textbf{Símbolo} & \textbf{Descripción} & \textbf{Unidad} \\
\midrule
$\kappa$                & \descKappa      & \unidadUSDt \\
$CAP_s$                 & \descCAPs       & tCO$_2$e \\
$CAP_d$                 & \descCAPd       & tCO$_2$e \\
$\sigma_\varepsilon^2$  & \descSigmaEps   & tCO$_2$e \\
$\sigma_{CAP}^2$        & \descSigmaCAP   & tCO$_2$e \\
\bottomrule
\end{tabular}}

% ---------------------------------------------------------------------
% 6. Valores base de los parámetros, con su origen
% ---------------------------------------------------------------------
% Las descripciones de origen se definen aparte para que la
% presentación pueda usarlas en tarjetas.
\newcommand{\origenCAPs}{\refcita{ministerio_de_energia_de_chile_plan_2024}}
\newcommand{\origenCAPd}{Promedio sector energía 2012--2022 (SNICHILE)}
\newcommand{\origenSigmaCAP}{Incertidumbre estructural histórica}
\newcommand{\origenKappa}{Costo social del carbono, Gobierno de Chile (2024)}

\newcommand{\tablaParametrosChile}{%
\begin{tabular}{l c l l}
\toprule
\textbf{Parámetro} & \textbf{Valor} & \textbf{Unidad} & \textbf{Origen} \\
\midrule
$CAP_s$          & \valCAPs      & \unidadMt   & \origenCAPs \\
$CAP_d$          & \valCAPd      & \unidadMt   & \origenCAPd \\
$\sigma_{CAP}$   & \valSigmaCAP  & \unidadMt   & \origenSigmaCAP \\
$\kappa$         & \valKappa     & \unidadUSDt & \origenKappa \\
\bottomrule
\end{tabular}}

% o
%   \newcommand{\refcita}[1]{\citet{#1}}   % presentación (natbib)
%   % =====================================================================
% shared/datos-modelo.tex
% ---------------------------------------------------------------------
% Nomenclatura, ecuaciones y valores de parámetros del modelo de
% anclaje y ajuste del subastador. Fuente única para la memoria
% (docs/MemoriaMHurtado) y la presentación de defensa
% (docs/Presentations/presentation-template), de modo que un valor o una
% ecuación se escriban una sola vez.
%
% Uso:
%   \newcommand{\refcita}[1]{\citeA{#1}}   % memoria (newapa/apacite)
%   \input{../../shared/datos-modelo}
% o
%   \newcommand{\refcita}[1]{\citet{#1}}   % presentación (natbib)
%   \input{../../../shared/datos-modelo}
%
% Los macros de tabla entregan sólo el `tabular`, sin flotante ni
% caption, para que cada documento lo envuelva a su manera (`table` en
% la memoria, `frame` en beamer).
%
% Requiere: amsmath, booktabs.
% =====================================================================

\providecommand{\refcita}[1]{\texttt{#1}}

% ---------------------------------------------------------------------
% 1. Valores de los parámetros calibrados para Chile
% ---------------------------------------------------------------------
% Valores tal como los reporta el cuadro de parámetros base de la
% memoria. \valCAPdExacto es el valor sin redondear que reproducen las
% tablas de resultados; el capítulo de conclusiones usa 311,01.
% Pendiente de unificar por los autores: ver
% docs/auditoria/AUDITORIA_2_MHurtado.md.
\newcommand{\valCAPs}{264,20}          % MtCO2e
\newcommand{\valCAPd}{311,00}          % MtCO2e
\newcommand{\valCAPdExacto}{311,0134}  % MtCO2e
\newcommand{\valSigmaCAP}{61,86}       % MtCO2e
\newcommand{\valKappa}{64,40}          % USD/tCO2e
% Calibración: hay que distinguir dos valores de m que la memoria usa para
% cosas distintas (chapter04.tex, Cuadro tab:top10_distancia_mix y el texto
% que le sigue).
%   - El argmin de la distancia absoluta es m=0, que exige sigma_eps -> infinito
%     y por eso no sirve como parámetro operativo.
%   - Para los análisis posteriores se selecciona m=0,1, el menor m con un
%     sigma_eps finito, cuya distancia (0,2758) queda muy cerca del mínimo.
% Nota: el texto de la memoria dice sigma_eps = 186,90 y el cuadro dice 185,58
% para el mismo m=0,1. Pendiente de unificar por los autores; ver
% docs/auditoria/AUDITORIA_2_MHurtado.md.
\newcommand{\valMArgmin}{0}                  % menor distancia absoluta
\newcommand{\valDistanciaArgmin}{0,2727}
\newcommand{\valMSeleccionado}{0,1}          % usado en los análisis posteriores
\newcommand{\valDistanciaSeleccionado}{0,2758}
\newcommand{\valSigmaEpsSeleccionado}{185,58} % MtCO2e, segun tab:top10_distancia_mix
\newcommand{\valHorizonte}{2020--2030}
\newcommand{\unidadMt}{MtCO$_2$e}
\newcommand{\unidadUSDt}{USD/tCO$_2$e}

% ---------------------------------------------------------------------
% 2. Ecuaciones del planificador con inatención conductual
% ---------------------------------------------------------------------
% Cuerpo matemático sin entorno, para poder usarlo en `equation`,
% en \[...\] o dentro de un nodo tikz.
\newcommand{\ecSenal}{CAP_s = CAP + \epsilon}
\newcommand{\ecEsperanza}{\mathbb{E}\left[CAP|CAP_s\right] = m CAP_s + (1 - m)CAP_d}
\newcommand{\ecAtencion}{m = \frac{\sigma_{CAP}^2}{\sigma_{CAP}^2 + \sigma_\varepsilon^2}}
\newcommand{\ecThetaSinPrecio}{{\theta^*}(CAP_s) = m\cdot CAP_s + (1 - m)\cdot CAP_d}
\newcommand{\ecUtilidadPlanificador}{\mathcal{V}(\theta,CAP) = \pi^a \theta - \frac{\kappa}{2}(\theta - CAP)^2}
\newcommand{\ecThetaOptimo}{\theta^*(CAP_s) = m\,CAP_s + (1-m)\,CAP_d + \frac{\pi^a}{\kappa}}

% ---------------------------------------------------------------------
% 3. Descripciones de los símbolos, definidas una sola vez
% ---------------------------------------------------------------------
% La memoria las usa como filas de tabla y la presentación como
% tarjetas; por eso el texto se define aquí y el formato en cada
% documento.
\newcommand{\descCAP}{Presupuesto ambiental real}
\newcommand{\descCAPs}{Señal observada del presupuesto de emisiones}
\newcommand{\descCAPd}{Presupuesto determinista del CAP (valor por defecto)}
\newcommand{\descEps}{Ruido en la señal}
\newcommand{\descSigmaEps}{Varianza del error en la señal}
\newcommand{\descSigmaCAP}{Varianza del CAP}
\newcommand{\descTheta}{Permisos totales emitidos por el planificador}
\newcommand{\descPiA}{Precio de permisos emitidos por el planificador}
\newcommand{\descPiV}{Precio de permisos en el mercado entre productores}
\newcommand{\descPiD}{Precio de la electricidad pagado por el consumidor}
\newcommand{\descKappa}{Costo social por alejarse del $CAP$}
\newcommand{\descM}{Nivel de atención del planificador}
% Distribuciones
\newcommand{\distCAP}{CAP \sim \mathcal{N}(CAP_d,\sigma^2_{CAP})}
\newcommand{\distEps}{\epsilon \sim \mathcal{N}(0,\sigma^2_\varepsilon)}

% ---------------------------------------------------------------------
% 4. Nomenclatura del modelo de Amigo et al. (2021)
% ---------------------------------------------------------------------
\newcommand{\tablaNomenclaturaAmigo}{%
\begin{tabular}{lll}
\toprule
\textbf{Símbolo} & \textbf{Descripción} & \textbf{Unidad} \\
\midrule
$\pi^{a}$ & \descPiA & USD/$tCO_2e$ \\
$\pi^{v}$ & \descPiV & USD/permiso \\
$\pi^{d}$ & \descPiD & USD/MWh \\
\midrule
$x_i$ & Inversión o expansión de capacidad del productor $i$ & MW \\
$Q_i$ & Generación eléctrica del productor $i$ & MWh \\
$A_i$ & Permisos adquiridos por el productor $i$ al planificador & $tCO_2e$\\
$P_i$ & Permisos comprados por el productor $i$ en el mercado secundario & $tCO_2e$ \\
$V_i$ & Permisos vendidos por el productor $i$ en el mercado secundario & $tCO_2e$ \\
$\epsilon_i$ & Factor de emisión por tecnología $i$ & $tCO_2$e/MWh \\
\midrule
$\theta$ & \descTheta & $tCO_2e$ \\
\bottomrule
\end{tabular}}

% ---------------------------------------------------------------------
% 5. Parámetros que el modelo conductual agrega a Amigo et al. (2021)
% ---------------------------------------------------------------------
\newcommand{\tablaParametrosAgregados}{%
\begin{tabular}{lll}
\toprule
\textbf{Símbolo} & \textbf{Descripción} & \textbf{Unidad} \\
\midrule
$\kappa$                & \descKappa      & \unidadUSDt \\
$CAP_s$                 & \descCAPs       & tCO$_2$e \\
$CAP_d$                 & \descCAPd       & tCO$_2$e \\
$\sigma_\varepsilon^2$  & \descSigmaEps   & tCO$_2$e \\
$\sigma_{CAP}^2$        & \descSigmaCAP   & tCO$_2$e \\
\bottomrule
\end{tabular}}

% ---------------------------------------------------------------------
% 6. Valores base de los parámetros, con su origen
% ---------------------------------------------------------------------
% Las descripciones de origen se definen aparte para que la
% presentación pueda usarlas en tarjetas.
\newcommand{\origenCAPs}{\refcita{ministerio_de_energia_de_chile_plan_2024}}
\newcommand{\origenCAPd}{Promedio sector energía 2012--2022 (SNICHILE)}
\newcommand{\origenSigmaCAP}{Incertidumbre estructural histórica}
\newcommand{\origenKappa}{Costo social del carbono, Gobierno de Chile (2024)}

\newcommand{\tablaParametrosChile}{%
\begin{tabular}{l c l l}
\toprule
\textbf{Parámetro} & \textbf{Valor} & \textbf{Unidad} & \textbf{Origen} \\
\midrule
$CAP_s$          & \valCAPs      & \unidadMt   & \origenCAPs \\
$CAP_d$          & \valCAPd      & \unidadMt   & \origenCAPd \\
$\sigma_{CAP}$   & \valSigmaCAP  & \unidadMt   & \origenSigmaCAP \\
$\kappa$         & \valKappa     & \unidadUSDt & \origenKappa \\
\bottomrule
\end{tabular}}

%
% Los macros de tabla entregan sólo el `tabular`, sin flotante ni
% caption, para que cada documento lo envuelva a su manera (`table` en
% la memoria, `frame` en beamer).
%
% Requiere: amsmath, booktabs.
% =====================================================================

\providecommand{\refcita}[1]{\texttt{#1}}

% ---------------------------------------------------------------------
% 1. Valores de los parámetros calibrados para Chile
% ---------------------------------------------------------------------
% Valores tal como los reporta el cuadro de parámetros base de la
% memoria. \valCAPdExacto es el valor sin redondear que reproducen las
% tablas de resultados; el capítulo de conclusiones usa 311,01.
% Pendiente de unificar por los autores: ver
% docs/auditoria/AUDITORIA_2_MHurtado.md.
\newcommand{\valCAPs}{264,20}          % MtCO2e
\newcommand{\valCAPd}{311,00}          % MtCO2e
\newcommand{\valCAPdExacto}{311,0134}  % MtCO2e
\newcommand{\valSigmaCAP}{61,86}       % MtCO2e
\newcommand{\valKappa}{64,40}          % USD/tCO2e
% Calibración: hay que distinguir dos valores de m que la memoria usa para
% cosas distintas (chapter04.tex, Cuadro tab:top10_distancia_mix y el texto
% que le sigue).
%   - El argmin de la distancia absoluta es m=0, que exige sigma_eps -> infinito
%     y por eso no sirve como parámetro operativo.
%   - Para los análisis posteriores se selecciona m=0,1, el menor m con un
%     sigma_eps finito, cuya distancia (0,2758) queda muy cerca del mínimo.
% Nota: el texto de la memoria dice sigma_eps = 186,90 y el cuadro dice 185,58
% para el mismo m=0,1. Pendiente de unificar por los autores; ver
% docs/auditoria/AUDITORIA_2_MHurtado.md.
\newcommand{\valMArgmin}{0}                  % menor distancia absoluta
\newcommand{\valDistanciaArgmin}{0,2727}
\newcommand{\valMSeleccionado}{0,1}          % usado en los análisis posteriores
\newcommand{\valDistanciaSeleccionado}{0,2758}
\newcommand{\valSigmaEpsSeleccionado}{185,58} % MtCO2e, segun tab:top10_distancia_mix
\newcommand{\valHorizonte}{2020--2030}
\newcommand{\unidadMt}{MtCO$_2$e}
\newcommand{\unidadUSDt}{USD/tCO$_2$e}

% ---------------------------------------------------------------------
% 2. Ecuaciones del planificador con inatención conductual
% ---------------------------------------------------------------------
% Cuerpo matemático sin entorno, para poder usarlo en `equation`,
% en \[...\] o dentro de un nodo tikz.
\newcommand{\ecSenal}{CAP_s = CAP + \epsilon}
\newcommand{\ecEsperanza}{\mathbb{E}\left[CAP|CAP_s\right] = m CAP_s + (1 - m)CAP_d}
\newcommand{\ecAtencion}{m = \frac{\sigma_{CAP}^2}{\sigma_{CAP}^2 + \sigma_\varepsilon^2}}
\newcommand{\ecThetaSinPrecio}{{\theta^*}(CAP_s) = m\cdot CAP_s + (1 - m)\cdot CAP_d}
\newcommand{\ecUtilidadPlanificador}{\mathcal{V}(\theta,CAP) = \pi^a \theta - \frac{\kappa}{2}(\theta - CAP)^2}
\newcommand{\ecThetaOptimo}{\theta^*(CAP_s) = m\,CAP_s + (1-m)\,CAP_d + \frac{\pi^a}{\kappa}}

% ---------------------------------------------------------------------
% 3. Descripciones de los símbolos, definidas una sola vez
% ---------------------------------------------------------------------
% La memoria las usa como filas de tabla y la presentación como
% tarjetas; por eso el texto se define aquí y el formato en cada
% documento.
\newcommand{\descCAP}{Presupuesto ambiental real}
\newcommand{\descCAPs}{Señal observada del presupuesto de emisiones}
\newcommand{\descCAPd}{Presupuesto determinista del CAP (valor por defecto)}
\newcommand{\descEps}{Ruido en la señal}
\newcommand{\descSigmaEps}{Varianza del error en la señal}
\newcommand{\descSigmaCAP}{Varianza del CAP}
\newcommand{\descTheta}{Permisos totales emitidos por el planificador}
\newcommand{\descPiA}{Precio de permisos emitidos por el planificador}
\newcommand{\descPiV}{Precio de permisos en el mercado entre productores}
\newcommand{\descPiD}{Precio de la electricidad pagado por el consumidor}
\newcommand{\descKappa}{Costo social por alejarse del $CAP$}
\newcommand{\descM}{Nivel de atención del planificador}
% Distribuciones
\newcommand{\distCAP}{CAP \sim \mathcal{N}(CAP_d,\sigma^2_{CAP})}
\newcommand{\distEps}{\epsilon \sim \mathcal{N}(0,\sigma^2_\varepsilon)}

% ---------------------------------------------------------------------
% 4. Nomenclatura del modelo de Amigo et al. (2021)
% ---------------------------------------------------------------------
\newcommand{\tablaNomenclaturaAmigo}{%
\begin{tabular}{lll}
\toprule
\textbf{Símbolo} & \textbf{Descripción} & \textbf{Unidad} \\
\midrule
$\pi^{a}$ & \descPiA & USD/$tCO_2e$ \\
$\pi^{v}$ & \descPiV & USD/permiso \\
$\pi^{d}$ & \descPiD & USD/MWh \\
\midrule
$x_i$ & Inversión o expansión de capacidad del productor $i$ & MW \\
$Q_i$ & Generación eléctrica del productor $i$ & MWh \\
$A_i$ & Permisos adquiridos por el productor $i$ al planificador & $tCO_2e$\\
$P_i$ & Permisos comprados por el productor $i$ en el mercado secundario & $tCO_2e$ \\
$V_i$ & Permisos vendidos por el productor $i$ en el mercado secundario & $tCO_2e$ \\
$\epsilon_i$ & Factor de emisión por tecnología $i$ & $tCO_2$e/MWh \\
\midrule
$\theta$ & \descTheta & $tCO_2e$ \\
\bottomrule
\end{tabular}}

% ---------------------------------------------------------------------
% 5. Parámetros que el modelo conductual agrega a Amigo et al. (2021)
% ---------------------------------------------------------------------
\newcommand{\tablaParametrosAgregados}{%
\begin{tabular}{lll}
\toprule
\textbf{Símbolo} & \textbf{Descripción} & \textbf{Unidad} \\
\midrule
$\kappa$                & \descKappa      & \unidadUSDt \\
$CAP_s$                 & \descCAPs       & tCO$_2$e \\
$CAP_d$                 & \descCAPd       & tCO$_2$e \\
$\sigma_\varepsilon^2$  & \descSigmaEps   & tCO$_2$e \\
$\sigma_{CAP}^2$        & \descSigmaCAP   & tCO$_2$e \\
\bottomrule
\end{tabular}}

% ---------------------------------------------------------------------
% 6. Valores base de los parámetros, con su origen
% ---------------------------------------------------------------------
% Las descripciones de origen se definen aparte para que la
% presentación pueda usarlas en tarjetas.
\newcommand{\origenCAPs}{\refcita{ministerio_de_energia_de_chile_plan_2024}}
\newcommand{\origenCAPd}{Promedio sector energía 2012--2022 (SNICHILE)}
\newcommand{\origenSigmaCAP}{Incertidumbre estructural histórica}
\newcommand{\origenKappa}{Costo social del carbono, Gobierno de Chile (2024)}

\newcommand{\tablaParametrosChile}{%
\begin{tabular}{l c l l}
\toprule
\textbf{Parámetro} & \textbf{Valor} & \textbf{Unidad} & \textbf{Origen} \\
\midrule
$CAP_s$          & \valCAPs      & \unidadMt   & \origenCAPs \\
$CAP_d$          & \valCAPd      & \unidadMt   & \origenCAPd \\
$\sigma_{CAP}$   & \valSigmaCAP  & \unidadMt   & \origenSigmaCAP \\
$\kappa$         & \valKappa     & \unidadUSDt & \origenKappa \\
\bottomrule
\end{tabular}}

%
% Los macros de tabla entregan sólo el `tabular`, sin flotante ni
% caption, para que cada documento lo envuelva a su manera (`table` en
% la memoria, `frame` en beamer).
%
% Requiere: amsmath, booktabs.
% =====================================================================

\providecommand{\refcita}[1]{\texttt{#1}}

% ---------------------------------------------------------------------
% 1. Valores de los parámetros calibrados para Chile
% ---------------------------------------------------------------------
% Valores tal como los reporta el cuadro de parámetros base de la
% memoria. \valCAPdExacto es el valor sin redondear que reproducen las
% tablas de resultados; el capítulo de conclusiones usa 311,01.
% Pendiente de unificar por los autores: ver
% docs/auditoria/AUDITORIA_2_MHurtado.md.
\newcommand{\valCAPs}{264,20}          % MtCO2e
\newcommand{\valCAPd}{311,00}          % MtCO2e
\newcommand{\valCAPdExacto}{311,0134}  % MtCO2e
\newcommand{\valSigmaCAP}{61,86}       % MtCO2e
\newcommand{\valKappa}{64,40}          % USD/tCO2e
% Calibración: hay que distinguir dos valores de m que la memoria usa para
% cosas distintas (chapter04.tex, Cuadro tab:top10_distancia_mix y el texto
% que le sigue).
%   - El argmin de la distancia absoluta es m=0, que exige sigma_eps -> infinito
%     y por eso no sirve como parámetro operativo.
%   - Para los análisis posteriores se selecciona m=0,1, el menor m con un
%     sigma_eps finito, cuya distancia (0,2758) queda muy cerca del mínimo.
% Nota: el texto de la memoria dice sigma_eps = 186,90 y el cuadro dice 185,58
% para el mismo m=0,1. Pendiente de unificar por los autores; ver
% docs/auditoria/AUDITORIA_2_MHurtado.md.
\newcommand{\valMArgmin}{0}                  % menor distancia absoluta
\newcommand{\valDistanciaArgmin}{0,2727}
\newcommand{\valMSeleccionado}{0,1}          % usado en los análisis posteriores
\newcommand{\valDistanciaSeleccionado}{0,2758}
\newcommand{\valSigmaEpsSeleccionado}{185,58} % MtCO2e, segun tab:top10_distancia_mix
\newcommand{\valHorizonte}{2020--2030}
\newcommand{\unidadMt}{MtCO$_2$e}
\newcommand{\unidadUSDt}{USD/tCO$_2$e}

% ---------------------------------------------------------------------
% 2. Ecuaciones del planificador con inatención conductual
% ---------------------------------------------------------------------
% Cuerpo matemático sin entorno, para poder usarlo en `equation`,
% en \[...\] o dentro de un nodo tikz.
\newcommand{\ecSenal}{CAP_s = CAP + \epsilon}
\newcommand{\ecEsperanza}{\mathbb{E}\left[CAP|CAP_s\right] = m CAP_s + (1 - m)CAP_d}
\newcommand{\ecAtencion}{m = \frac{\sigma_{CAP}^2}{\sigma_{CAP}^2 + \sigma_\varepsilon^2}}
\newcommand{\ecThetaSinPrecio}{{\theta^*}(CAP_s) = m\cdot CAP_s + (1 - m)\cdot CAP_d}
\newcommand{\ecUtilidadPlanificador}{\mathcal{V}(\theta,CAP) = \pi^a \theta - \frac{\kappa}{2}(\theta - CAP)^2}
\newcommand{\ecThetaOptimo}{\theta^*(CAP_s) = m\,CAP_s + (1-m)\,CAP_d + \frac{\pi^a}{\kappa}}

% ---------------------------------------------------------------------
% 3. Descripciones de los símbolos, definidas una sola vez
% ---------------------------------------------------------------------
% La memoria las usa como filas de tabla y la presentación como
% tarjetas; por eso el texto se define aquí y el formato en cada
% documento.
\newcommand{\descCAP}{Presupuesto ambiental real}
\newcommand{\descCAPs}{Señal observada del presupuesto de emisiones}
\newcommand{\descCAPd}{Presupuesto determinista del CAP (valor por defecto)}
\newcommand{\descEps}{Ruido en la señal}
\newcommand{\descSigmaEps}{Varianza del error en la señal}
\newcommand{\descSigmaCAP}{Varianza del CAP}
\newcommand{\descTheta}{Permisos totales emitidos por el planificador}
\newcommand{\descPiA}{Precio de permisos emitidos por el planificador}
\newcommand{\descPiV}{Precio de permisos en el mercado entre productores}
\newcommand{\descPiD}{Precio de la electricidad pagado por el consumidor}
\newcommand{\descKappa}{Costo social por alejarse del $CAP$}
\newcommand{\descM}{Nivel de atención del planificador}
% Distribuciones
\newcommand{\distCAP}{CAP \sim \mathcal{N}(CAP_d,\sigma^2_{CAP})}
\newcommand{\distEps}{\epsilon \sim \mathcal{N}(0,\sigma^2_\varepsilon)}

% ---------------------------------------------------------------------
% 4. Nomenclatura del modelo de Amigo et al. (2021)
% ---------------------------------------------------------------------
\newcommand{\tablaNomenclaturaAmigo}{%
\begin{tabular}{lll}
\toprule
\textbf{Símbolo} & \textbf{Descripción} & \textbf{Unidad} \\
\midrule
$\pi^{a}$ & \descPiA & USD/$tCO_2e$ \\
$\pi^{v}$ & \descPiV & USD/permiso \\
$\pi^{d}$ & \descPiD & USD/MWh \\
\midrule
$x_i$ & Inversión o expansión de capacidad del productor $i$ & MW \\
$Q_i$ & Generación eléctrica del productor $i$ & MWh \\
$A_i$ & Permisos adquiridos por el productor $i$ al planificador & $tCO_2e$\\
$P_i$ & Permisos comprados por el productor $i$ en el mercado secundario & $tCO_2e$ \\
$V_i$ & Permisos vendidos por el productor $i$ en el mercado secundario & $tCO_2e$ \\
$\epsilon_i$ & Factor de emisión por tecnología $i$ & $tCO_2$e/MWh \\
\midrule
$\theta$ & \descTheta & $tCO_2e$ \\
\bottomrule
\end{tabular}}

% ---------------------------------------------------------------------
% 5. Parámetros que el modelo conductual agrega a Amigo et al. (2021)
% ---------------------------------------------------------------------
\newcommand{\tablaParametrosAgregados}{%
\begin{tabular}{lll}
\toprule
\textbf{Símbolo} & \textbf{Descripción} & \textbf{Unidad} \\
\midrule
$\kappa$                & \descKappa      & \unidadUSDt \\
$CAP_s$                 & \descCAPs       & tCO$_2$e \\
$CAP_d$                 & \descCAPd       & tCO$_2$e \\
$\sigma_\varepsilon^2$  & \descSigmaEps   & tCO$_2$e \\
$\sigma_{CAP}^2$        & \descSigmaCAP   & tCO$_2$e \\
\bottomrule
\end{tabular}}

% ---------------------------------------------------------------------
% 6. Valores base de los parámetros, con su origen
% ---------------------------------------------------------------------
% Las descripciones de origen se definen aparte para que la
% presentación pueda usarlas en tarjetas.
\newcommand{\origenCAPs}{\refcita{ministerio_de_energia_de_chile_plan_2024}}
\newcommand{\origenCAPd}{Promedio sector energía 2012--2022 (SNICHILE)}
\newcommand{\origenSigmaCAP}{Incertidumbre estructural histórica}
\newcommand{\origenKappa}{Costo social del carbono, Gobierno de Chile (2024)}

\newcommand{\tablaParametrosChile}{%
\begin{tabular}{l c l l}
\toprule
\textbf{Parámetro} & \textbf{Valor} & \textbf{Unidad} & \textbf{Origen} \\
\midrule
$CAP_s$          & \valCAPs      & \unidadMt   & \origenCAPs \\
$CAP_d$          & \valCAPd      & \unidadMt   & \origenCAPd \\
$\sigma_{CAP}$   & \valSigmaCAP  & \unidadMt   & \origenSigmaCAP \\
$\kappa$         & \valKappa     & \unidadUSDt & \origenKappa \\
\bottomrule
\end{tabular}}

%
% Los macros de tabla entregan sólo el `tabular`, sin flotante ni
% caption, para que cada documento lo envuelva a su manera (`table` en
% la memoria, `frame` en beamer).
%
% Requiere: amsmath, booktabs.
% =====================================================================

\providecommand{\refcita}[1]{\texttt{#1}}

% ---------------------------------------------------------------------
% 1. Valores de los parámetros calibrados para Chile
% ---------------------------------------------------------------------
% Valores tal como los reporta el cuadro de parámetros base de la
% memoria. \valCAPdExacto es el valor sin redondear que reproducen las
% tablas de resultados; el capítulo de conclusiones usa 311,01.
% Pendiente de unificar por los autores: ver
% docs/auditoria/AUDITORIA_2_MHurtado.md.
\newcommand{\valCAPs}{264,20}          % MtCO2e
\newcommand{\valCAPd}{311,00}          % MtCO2e
\newcommand{\valCAPdExacto}{311,0134}  % MtCO2e
\newcommand{\valSigmaCAP}{61,86}       % MtCO2e
\newcommand{\valKappa}{64,40}          % USD/tCO2e
% Calibración: hay que distinguir dos valores de m que la memoria usa para
% cosas distintas (chapter04.tex, Cuadro tab:top10_distancia_mix y el texto
% que le sigue).
%   - El argmin de la distancia absoluta es m=0, que exige sigma_eps -> infinito
%     y por eso no sirve como parámetro operativo.
%   - Para los análisis posteriores se selecciona m=0,1, el menor m con un
%     sigma_eps finito, cuya distancia (0,2758) queda muy cerca del mínimo.
% Nota: el texto de la memoria dice sigma_eps = 186,90 y el cuadro dice 185,58
% para el mismo m=0,1. Pendiente de unificar por los autores; ver
% docs/auditoria/AUDITORIA_2_MHurtado.md.
\newcommand{\valMArgmin}{0}                  % menor distancia absoluta
\newcommand{\valDistanciaArgmin}{0,2727}
\newcommand{\valMSeleccionado}{0,1}          % usado en los análisis posteriores
\newcommand{\valDistanciaSeleccionado}{0,2758}
\newcommand{\valSigmaEpsSeleccionado}{185,58} % MtCO2e, segun tab:top10_distancia_mix
\newcommand{\valHorizonte}{2020--2030}
\newcommand{\unidadMt}{MtCO$_2$e}
\newcommand{\unidadUSDt}{USD/tCO$_2$e}

% ---------------------------------------------------------------------
% 2. Ecuaciones del planificador con inatención conductual
% ---------------------------------------------------------------------
% Cuerpo matemático sin entorno, para poder usarlo en `equation`,
% en \[...\] o dentro de un nodo tikz.
\newcommand{\ecSenal}{CAP_s = CAP + \epsilon}
\newcommand{\ecEsperanza}{\mathbb{E}\left[CAP|CAP_s\right] = m CAP_s + (1 - m)CAP_d}
\newcommand{\ecAtencion}{m = \frac{\sigma_{CAP}^2}{\sigma_{CAP}^2 + \sigma_\varepsilon^2}}
\newcommand{\ecThetaSinPrecio}{{\theta^*}(CAP_s) = m\cdot CAP_s + (1 - m)\cdot CAP_d}
\newcommand{\ecUtilidadPlanificador}{\mathcal{V}(\theta,CAP) = \pi^a \theta - \frac{\kappa}{2}(\theta - CAP)^2}
\newcommand{\ecThetaOptimo}{\theta^*(CAP_s) = m\,CAP_s + (1-m)\,CAP_d + \frac{\pi^a}{\kappa}}

% ---------------------------------------------------------------------
% 3. Descripciones de los símbolos, definidas una sola vez
% ---------------------------------------------------------------------
% La memoria las usa como filas de tabla y la presentación como
% tarjetas; por eso el texto se define aquí y el formato en cada
% documento.
\newcommand{\descCAP}{Presupuesto ambiental real}
\newcommand{\descCAPs}{Señal observada del presupuesto de emisiones}
\newcommand{\descCAPd}{Presupuesto determinista del CAP (valor por defecto)}
\newcommand{\descEps}{Ruido en la señal}
\newcommand{\descSigmaEps}{Varianza del error en la señal}
\newcommand{\descSigmaCAP}{Varianza del CAP}
\newcommand{\descTheta}{Permisos totales emitidos por el planificador}
\newcommand{\descPiA}{Precio de permisos emitidos por el planificador}
\newcommand{\descPiV}{Precio de permisos en el mercado entre productores}
\newcommand{\descPiD}{Precio de la electricidad pagado por el consumidor}
\newcommand{\descKappa}{Costo social por alejarse del $CAP$}
\newcommand{\descM}{Nivel de atención del planificador}
% Distribuciones
\newcommand{\distCAP}{CAP \sim \mathcal{N}(CAP_d,\sigma^2_{CAP})}
\newcommand{\distEps}{\epsilon \sim \mathcal{N}(0,\sigma^2_\varepsilon)}

% ---------------------------------------------------------------------
% 4. Nomenclatura del modelo de Amigo et al. (2021)
% ---------------------------------------------------------------------
\newcommand{\tablaNomenclaturaAmigo}{%
\begin{tabular}{lll}
\toprule
\textbf{Símbolo} & \textbf{Descripción} & \textbf{Unidad} \\
\midrule
$\pi^{a}$ & \descPiA & USD/$tCO_2e$ \\
$\pi^{v}$ & \descPiV & USD/permiso \\
$\pi^{d}$ & \descPiD & USD/MWh \\
\midrule
$x_i$ & Inversión o expansión de capacidad del productor $i$ & MW \\
$Q_i$ & Generación eléctrica del productor $i$ & MWh \\
$A_i$ & Permisos adquiridos por el productor $i$ al planificador & $tCO_2e$\\
$P_i$ & Permisos comprados por el productor $i$ en el mercado secundario & $tCO_2e$ \\
$V_i$ & Permisos vendidos por el productor $i$ en el mercado secundario & $tCO_2e$ \\
$\epsilon_i$ & Factor de emisión por tecnología $i$ & $tCO_2$e/MWh \\
\midrule
$\theta$ & \descTheta & $tCO_2e$ \\
\bottomrule
\end{tabular}}

% ---------------------------------------------------------------------
% 5. Parámetros que el modelo conductual agrega a Amigo et al. (2021)
% ---------------------------------------------------------------------
\newcommand{\tablaParametrosAgregados}{%
\begin{tabular}{lll}
\toprule
\textbf{Símbolo} & \textbf{Descripción} & \textbf{Unidad} \\
\midrule
$\kappa$                & \descKappa      & \unidadUSDt \\
$CAP_s$                 & \descCAPs       & tCO$_2$e \\
$CAP_d$                 & \descCAPd       & tCO$_2$e \\
$\sigma_\varepsilon^2$  & \descSigmaEps   & tCO$_2$e \\
$\sigma_{CAP}^2$        & \descSigmaCAP   & tCO$_2$e \\
\bottomrule
\end{tabular}}

% ---------------------------------------------------------------------
% 6. Valores base de los parámetros, con su origen
% ---------------------------------------------------------------------
% Las descripciones de origen se definen aparte para que la
% presentación pueda usarlas en tarjetas.
\newcommand{\origenCAPs}{\refcita{ministerio_de_energia_de_chile_plan_2024}}
\newcommand{\origenCAPd}{Promedio sector energía 2012--2022 (SNICHILE)}
\newcommand{\origenSigmaCAP}{Incertidumbre estructural histórica}
\newcommand{\origenKappa}{Costo social del carbono, Gobierno de Chile (2024)}

\newcommand{\tablaParametrosChile}{%
\begin{tabular}{l c l l}
\toprule
\textbf{Parámetro} & \textbf{Valor} & \textbf{Unidad} & \textbf{Origen} \\
\midrule
$CAP_s$          & \valCAPs      & \unidadMt   & \origenCAPs \\
$CAP_d$          & \valCAPd      & \unidadMt   & \origenCAPd \\
$\sigma_{CAP}$   & \valSigmaCAP  & \unidadMt   & \origenSigmaCAP \\
$\kappa$         & \valKappa     & \unidadUSDt & \origenKappa \\
\bottomrule
\end{tabular}}
